\documentclass{amsart}
\usepackage{amssymb,mathtools,array,longtable}
\usepackage[hidelinks]{hyperref}
\setlength{\textwidth}{32pc}
\setlength{\textheight}{50pc}
\setlength{\oddsidemargin}{\dimexpr(\paperwidth-\textwidth)/2-1in\relax}
\setlength{\evensidemargin}{\oddsidemargin}
\setlength{\topmargin}{\dimexpr(\paperheight-\textheight-\headheight-\headsep)/2-1in\relax}
\emergencystretch=1.5em
\hypersetup{pdftitle={Irrationality of the dilogarithm at 1/3},pdfauthor={Ruiyi Zhang}}

\makeatletter
\def\LT@makecaption#1#2#3{\LT@mcol\LT@cols c{\hbox to\z@{\hss
  \parbox[t]\textwidth{\normalsize\@tempcnta=64 \@makecaption{#2}{#3}}\hss}}}
\makeatother

\newtheorem{theorem}{Theorem}
\newtheorem{lemma}{Lemma}[section]
\newtheorem{proposition}[lemma]{Proposition}
\newtheorem{fact}{Fact}
\numberwithin{equation}{section}
\numberwithin{table}{section}

\newcommand{\Li}{\operatorname{Li}}
\newcommand{\Z}{\mathbb{Z}}
\makeatletter
\newcommand{\captionof}[1]{\def\@captype{#1}\caption}
\makeatother

% ---------------------------------------------------------------------------
% Theorem-level numbers (generated from the certified data).
\newcommand{\Rseam}{8}
\newcommand{\nslabs}{76}
\newcommand{\innervalue}{\frac{580566971}{177365760}}
\newcommand{\Avalue}{\frac{902042411}{177365760}}
\newcommand{\tauvalue}{3.872724286099666640}
\newcommand{\taubound}{3.87273}
\newcommand{\qminus}{-5.193572}
\newcommand{\Kx}{-4.83148}
\newcommand{\Kv}{-4.74318}
\newcommand{\Kz}{0.04318}
\newcommand{\Kzcirc}{0.04283}
\newcommand{\alphazero}{3.875784}
\newcommand{\alphabound}{3.8757}
\newcommand{\marginvalue}{0.00305}
\newcommand{\capmargin}{2\cdot10^{-6}}
\newcommand{\Narcs}{32}
\newcommand{\nleaves}{126464}
\newcommand{\NE}{150}
\newcommand{\NF}{90}
\newcommand{\precE}{320}
\newcommand{\precF}{192}
\newcommand{\vzero}{0.00960735979838}
% Line-family index ranges for the seam R=\Rseam (integer R>=5):
%   k<R/4, k<R/2, k<3R/2, k<5R/2, w<R/4+1  <=>  k<=floor((R-1)/4), ...
\newcommand{\fdiv}[2]{\the\numexpr(2*(#1)-(#2)+1)/(2*(#2))\relax}% floor(a/b), a>=0
\newcommand{\kLlo}{\fdiv{\Rseam-1}{4}}
\newcommand{\kKlo}{\fdiv{\Rseam-1}{2}}
\newcommand{\kKhi}{\fdiv{3*\Rseam-1}{2}}
\newcommand{\kLhi}{\fdiv{5*\Rseam-1}{2}}
\newcommand{\wmax}{\fdiv{\Rseam+3}{4}}
\newcommand{\nlines}{\the\numexpr 2+(\kLlo+1)+(\kKlo+1)+\kKhi+(\kLhi-1)+\wmax\relax}
\newcommand{\Rfour}{\the\numexpr4*\Rseam\relax}
% ---------------------------------------------------------------------------

\begin{document}

\title{Irrationality of the dilogarithm at $1/3$}
\author{Ruiyi Zhang}
\address{Independent Researcher, Columbia, Maryland, USA}
\email{ryzfsu@gmail.com}
\subjclass[2020]{Primary 11J72; Secondary 33B30}
\keywords{dilogarithm, irrationality, interpolation determinant, logarithmic energy, interval arithmetic}

\begin{abstract}
We prove that $\operatorname{Li}_2(1/3)$ is irrational. Irrationality of
$\operatorname{Li}_2(1/s)$ was known for integers $s\le-5$ and $s\ge6$ by
work of Hata and of Rhin and Viola, and Prellberg recently proved it for
$s=-4,5,-3$. We define polynomials $G_m(T)$ with rational coefficients by an
interpolation determinant. The rank of this determinant modulo primes
$p$ with $m/\Rseam<p<5m/4$ bounds the denominators of the coefficients. A
concavity inequality for a logarithmic energy, combined with the inversion
$v\mapsto3/v$, bounds $G_m(\operatorname{Li}_2(1/3))$. Interval
arithmetic is used only for the energy and the potentials of a
printed reference measure, on a uniform dyadic subdivision. The integral of the local rank function is a rational
sum over $\nslabs$ printed affine pieces. The coefficients of $G_4$
and the $9\times9$ determinant modulo $17$ are printed and can be checked by hand. A
congruence modulo $17$ shows that $G_m(\xi)\ne0$ for every rational $\xi$
and $m=4\cdot17^t$, $t$ large in terms of $\xi$.
\end{abstract}

\maketitle

\section{Introduction}\label{sec:intro}

For $|z|\le1$ let $\Li_2(z)=\sum_{k\ge1}z^k/k^2$.
Hata \cite[Theorem~1.1]{Ha} proved that $\Li_2(1/s)$ is irrational
for every integer $s\le-5$ and every $s\ge7$.
Rhin and Viola settled $s=6$ by their permutation group method
\cite[Section~6]{RV05}, and later treated negative arguments
\cite[Theorem~2.1 and Section~4]{RV19}.
Viola and Zudilin \cite{VZ} introduced a six-parameter family of double
integrals for dilogarithmic values.
In 2024 Calegari, Dimitrov and Tang
\cite[Remark~2.8.2 and Section~15.2]{CDT} listed the integers still open:
$s\in\{-4,-3,-2,2,3,4,5\}$.
Prellberg's 2026 preprint \cite[Theorem~1.1]{Pr} proves the cases $s=-4,5,-3$.
We prove the case $s=3$.
Our linear forms are values of an interpolation determinant of size $9n$,
in place of the double integrals of \cite{Ha,Pr,RV05,VZ}. The rank of this
determinant modulo primes controls the denominators of its coefficients.

\begin{theorem}\label{thm:main}
The number $\Li_2(1/3)$ is irrational.
\end{theorem}

By Landen's identity \cite[(25.12.3)]{DLMF},
$\Li_2(1/3)+\Li_2(-1/2)=-\frac12\log^2(3/2)$, and $\log(3/2)$ is
transcendental by the Hermite--Lindemann theorem, a special case of
\cite[Theorem~1.4]{Baker}. So Landen's identity does not reduce the case
$s=3$ to the case $s=-2$.

Put $\theta=\Li_2(1/3)$ and $m=4n$. Section~\ref{sec:determinant} defines
$G_m(T)\in\mathbb Q[T]$ by an interpolation determinant.
Sections~\ref{sec:divisor} and~\ref{sec:analytic} give
\begin{gather}
\deg G_m\le m\ \text{by Lemma~\ref{lem:moment-identity}},\qquad
K_mG_m(T)\in\Z[T],
\qquad K_m\in\mathbb Q_{>0},\notag
\\
\limsup_{\substack{m\to\infty\\4\mid m}}
  \frac{\log K_m}{m^2}\le\tau<\taubound,\notag\\
\log|G_m(\theta)|\le-\alphabound\,m^2+C\,m\log m
\quad(m=4n\ge4).\label{eq:intro-rates}
\end{gather}
The exact value of $\tau$ is in \eqref{eq:tau}, and $\alphabound$ is a
lower bound for the constant of \eqref{eq:alpha}. The finite constant $C$ of
Lemma~\ref{lem:analytic-bound} depends only on the reference measure.
Since $\taubound<\alphabound$, the quadratic exponent is negative.
If $\theta=u/q_0$ with $q_0>0$, then $q_0^mK_mG_m(\theta)\in\Z$.
Section~\ref{sec:nonvanishing} shows that it is nonzero for
$m=4\cdot17^t$ and $t$ large, and \eqref{eq:intro-rates} shows that it tends
to zero.

The only numerical margin in the proof is $\alpha_0-\tau>0$. The size of
$C$ is not used.
Lemma~\ref{lem:reference-bound} is computer-assisted. It is proved by
interval arithmetic on the reference measure of
Tables~\ref{tab:reference-intervals} and~\ref{tab:reference-poisson}, over
the uniform subdivision of Lemma~\ref{lem:leaves}. The integral
\eqref{eq:inner-integral} is an exact sum over the $\nslabs$ affine pieces of
Table~\ref{tab:slabs}, each determined by the explicit lines of
Table~\ref{tab:line-families}.
Tables~\ref{tab:g4-coefficients}, \ref{tab:line-families} and~\ref{tab:slabs}
and \eqref{eq:g4-mod-matrix} can be checked by hand. Python scripts recheck
Lemma~\ref{lem:reference-bound} and part of the remaining proof
(Appendix~\ref{app:ver}). Apart from this lemma, the proof is complete without them.

The proof uses three classical inputs.

\begin{fact}[prime number theorem {\cite[Theorems~6 and~420, \S\S1.8 and~22.4]{HW}}]\label{fact:pnt}
$\theta_{\rm Ch}(y)=\sum_{p\le y}\log p$ satisfies $\theta_{\rm Ch}(y)/y\to1$
as $y\to\infty$.
\end{fact}
Only this qualitative form is used.

\begin{fact}[{\cite[Lemma~I.1.8]{ST}}]\label{fact:energy}
If $\sigma$ is a real compactly supported signed measure of total mass zero
with $\iint|\log|u-w||\,d|\sigma|(u)\,d|\sigma|(w)<\infty$, then
$\iint\log|u-w|\,d\sigma(u)\,d\sigma(w)\le0$.
\end{fact}
A short proof is given in Appendix~\ref{app:energy}.

\begin{fact}[Cauchy--Vandermonde {\cite[Theorem~25]{Krat}, \cite[Section~2]{HK}}]\label{fact:cv}
For $s=(s_1,\ldots,s_m)$, $v=(v_1,\ldots,v_L)$ and $h=L-m\ge0$,
\[
\det\begin{pmatrix}
((s_{i'}+v_b)^{-1})_{1\le i'\le m,\,1\le b\le L}\\(v_b^{\ell-1})_{1\le\ell\le h,\,1\le b\le L}
\end{pmatrix}
=\frac{\Delta(s)\Delta(v)}{\prod_{i',b}(s_{i'}+v_b)},
\qquad \Delta(u)=\prod_{i<j}(u_j-u_i).
\]
\end{fact}
A proof, including the sign, is given in the proof of
Lemma~\ref{lem:integral-identity}.

The use of AI tools is described at the end of the paper.

\section{The determinant}\label{sec:determinant}

Put $a=3$. For $n\ge1$ let
\begin{equation*}
\begin{gathered}
m=4n,\qquad \sigma=k=h=n,\qquad L=m+h,\\
I=\{\sigma+k,\ldots,\sigma+k+m-1\},\quad
J=\{\sigma,\ldots,\sigma+2m+h-1\},\quad
\Lambda=\{1,\ldots,h\},
\end{gathered}
\end{equation*}
all ordered increasingly. For $j\ge0$ put
\[
t_j=\sum_{u=1}^{j}\frac{a^{j-u}}u,\qquad
b_j=\sum_{u=1}^{j}\frac{a^{j-u}}{u^2};\qquad t_0=b_0=0.
\]
For $i\in I$, $j\in J$ and $\ell\in\Lambda$ let
\begin{equation*}
\begin{gathered}
C_{ij}=\begin{cases}(a^i-a^j)/(j-i),&i\ne j,\\0,&i=j,\end{cases}
\quad
B_{ij}=\begin{cases}(t_i-t_j)/(j-i),&i\ne j,\\b_i,&i=j,\end{cases}\\
A_{ij}=a^i\delta_{ij},\qquad V_{\ell j}=\frac1{j+\ell},\qquad
M_{ij}=\int_0^1\!\int_0^1\frac{x^iy^j}{xy-3}\,dx\,dy,
\end{gathered}
\end{equation*}
and put
\begin{equation}\label{eq:determinant}
G_m(T)=\det\begin{pmatrix}C\\B-TA\\V\end{pmatrix},
\end{equation}
a determinant of size $2m+h=9n$. For an indeterminate $a$ the same formulas
define
\[
G_m(T;a)=\det\begin{pmatrix}C(a)\\B(a)-TA(a)\\V\end{pmatrix}
\in\mathbb Q[a,T];\qquad G_m(T)=G_m(T;3).
\]

\begin{lemma}\label{lem:moment-identity}
One has $G_m\in\mathbb Q[T]$, $\deg G_m\le m$, and
\begin{equation}\label{eq:moment-identity}
M=B+\log(2/3)\,C-\theta A,\qquad
G_m(\theta)=\det\begin{pmatrix}C\\M\\V\end{pmatrix}.
\end{equation}
\end{lemma}

\begin{proof}
Write $\lambda=\log(2/3)$ and
\[
\mathfrak m_j=\int_0^1\frac{u^j}{u-3}\,du,\qquad
\mathfrak n_j=\int_0^1\frac{u^j\log(1/u)}{u-3}\,du.
\]
Polynomial division and the geometric series give
\begin{equation*}
\begin{gathered}
\mathfrak m_0=\lambda,\quad \mathfrak n_0=-\theta,\qquad
\mathfrak m_j=a\mathfrak m_{j-1}+\frac1j,\quad
\mathfrak n_j=a\mathfrak n_{j-1}+\frac1{j^2}\quad(j\ge1),\\
\mathfrak m_j=a^j\lambda+t_j,\qquad \mathfrak n_j=b_j-a^j\theta.
\end{gathered}
\end{equation*}
The substitution $u=xy$ gives
\[
M_{ij}=\frac{\mathfrak m_i-\mathfrak m_j}{j-i}\quad(i\ne j),\qquad M_{ii}=\mathfrak n_i.
\]
Hence the first identity in \eqref{eq:moment-identity}. Add $\lambda$
times each $C$ row to the corresponding row of $B-\theta A$. Only the $m$
middle rows depend on $T$, and their coefficients are rational.
\end{proof}

Put $r=2/5$ and $\gamma(\vartheta)=re^{i\vartheta}$ $(0\le\vartheta\le2\pi)$.
Let $[3,1]$ be the segment from $3$ to $1$, and let
\begin{equation*}
\Gamma=[3,1]+c_\Gamma\gamma,\qquad c_\Gamma=\frac{\log3}{2\pi i},
\end{equation*}
a linear combination of paths with complex coefficients. For every integer $d$,
\begin{equation}\label{eq:laurent-moments}
\int_\Gamma z^{d-1}\,dz=
\begin{cases}
(1-a^d)/d,&d\ne0,\\
0,&d=0.
\end{cases}
\end{equation}
For $d\ne0$ integrate $z^d/d$ on each path. For $d=0$ the segment
gives $-\log3$ and the circle gives $2\pi ic_\Gamma=\log3$. The total variation is
\begin{equation}\label{eq:chain-variation}
T_r=2+r\log3=2.43944\ldots<3.
\end{equation}

With $\Delta$ as in Fact~\ref{fact:cv}, use variables
$x=(x_1,\ldots,x_m)$, $v=(v_1,\ldots,v_L)$ and $z=(z_1,\ldots,z_m)$.

\begin{lemma}\label{lem:integral-identity}
With product differentials $dx\,dv\,dz$,
\begin{equation}\label{eq:integral-identity}
G_m(\theta)=
\frac{3^{3m^2/2-m}}{(m!)^2L!}
\int_{[0,1]^m}\int_{[0,1]^L}\int_{\Gamma^m}
 \Phi(x,v,z)\,dz\,dv\,dx,
\end{equation}
where
\begin{equation*}
\Phi(x,v,z)=
\frac{\Delta(x)^2\Delta(v)^2\Delta(z)^2
\displaystyle\prod_{d=1}^{m}\prod_{b=1}^{L}(z_d-v_b)
\displaystyle\prod_{i'=1}^{m}x_{i'}^{\sigma+k+h}
\prod_{b=1}^{L}v_b^\sigma\prod_{d=1}^{m}z_d^{-m-k}}
{\displaystyle\prod_{i'=1}^{m}\prod_{b=1}^{L}(3-x_{i'}v_b)}.
\end{equation*}
\end{lemma}

\begin{proof}
By \eqref{eq:laurent-moments},
$C_{ij}=a^i\int_\Gamma z^{j-i-1}\,dz$.
Put $f_i(v)=\int_0^1x^i/(xv-3)\,dx$. Then
\[
M_{ij}=\int_0^1f_i(v)v^j\,dv,\qquad
V_{\ell j}=\int_0^1v^{\ell-1}v^j\,dv.
\]
Set $e_\Delta=m(m-1)/2$ and
$\mathbf w=(z_1,\ldots,z_m,v_1,\ldots,v_L)$.
Multilinearity in the rows and antisymmetrization, applied to the $L$ rows
$[M;V]$, then to the $m$ rows $C$, and then to the $m$ functions $f_i$, give
\begin{equation*}
G_m(\theta)=\frac1{(m!)^2L!}
\int_{[0,1]^m\times[0,1]^L\times\Gamma^m}
D_zD_{zv}D_xD_{xv}\,dx\,dv\,dz,
\end{equation*}
where
\begin{align*}
D_z&=\det(a^iz_d^{-i-1})_{i\in I,\,1\le d\le m}
 =(-1)^{e_\Delta}a^{m(\sigma+k)+e_\Delta}\Delta(z)\prod_dz_d^{-\sigma-k-m},\\
D_{zv}&=\det(\mathbf w_b^j)_{1\le b\le 2m+h,\,j\in J}
=(-1)^{mL}\Delta(z)\Delta(v)
  \prod_dz_d^\sigma\prod_bv_b^\sigma\prod_{d,b}(z_d-v_b),\\
D_x&=\det(x_{i'}^i)_{i\in I,\,1\le i'\le m}
 =\Delta(x)\prod_{i'}x_{i'}^{\sigma+k},\\
D_{xv}&=\det\begin{pmatrix}
((x_{i'}v_b-3)^{-1})_{1\le i'\le m,\,1\le b\le L}\\
(v_b^{\ell-1})_{1\le\ell\le h,\,1\le b\le L}
\end{pmatrix}
 =\frac{(-1)^{mL}3^{e_\Delta}\Delta(x)\Delta(v)\prod_{i'}x_{i'}^h}
{\prod_{i',b}(3-x_{i'}v_b)}.
\end{align*}
The last equality follows from the Cauchy--Vandermonde identity. Explicitly,
for $s_{i'}=-3/x_{i'}$,
\[
\det\begin{pmatrix}((x_{i'}v_b-3)^{-1})_{i',b}\\(v_b^{\ell-1})_{\ell,b}\end{pmatrix}
=\prod_{i'}x_{i'}^{-1}
\det\begin{pmatrix}((s_{i'}+v_b)^{-1})_{i',b}\\(v_b^{\ell-1})_{\ell,b}\end{pmatrix},
\]
and $\Delta(s)=3^{e_\Delta}\Delta(x)\prod_{i'}x_{i'}^{1-m}$,
$\prod_{i',b}(s_{i'}+v_b)=(-1)^{mL}\prod_{i',b}(3-x_{i'}v_b)\prod_{i'}x_{i'}^{-L}$.
Fact~\ref{fact:cv} and this row scaling give the formula for $D_{xv}$.
Continuity covers $x_{i'}=0$.

We prove Fact~\ref{fact:cv} as an identity of rational
functions, with $s_{i'},v_b$ independent indeterminates. We use induction on
$m$, for all $h\ge0$ at once. The case $m=0$ is the Vandermonde identity.
Let $m\ge1$, let $v'=(v_1,\ldots,v_{L-1})$, and let $\widehat s_{i'}$
denote $s$ with $s_{i'}$ omitted. As a function of $v_L$, expansion along the
last column writes the determinant as
$\sum_{i'}(-1)^{i'+L}M_{i'}/(s_{i'}+v_L)$ plus a polynomial of degree at most
$h-1$, where $M_{i'}$ is the minor without row $i'$ and column $L$. The right
side is a polynomial of degree $L-1=m+h-1$ in $v_L$ divided by
$\prod_{i'}(s_{i'}+v_L)$, so it has the same form. By the induction
hypothesis for $(\widehat s_{i'},v')$, $M_{i'}$ is the right side of the
fact for $(\widehat s_{i'},v')$. We have
$\Delta(s)=(-1)^{i'-1}\Delta(\widehat s_{i'})\prod_{j\ne i'}(s_j-s_{i'})$,
and $\Delta(v)=(-1)^{L-1}\Delta(v')\prod_{b<L}(s_{i'}+v_b)$ at
$v_L=-s_{i'}$. Hence, after the
factors $s_j+v_L=s_j-s_{i'}$ $(j\ne i')$ and $s_{i'}+v_b$ $(b<L)$ cancel,
the residue of the right side at $v_L=-s_{i'}$ is
\[
 (-1)^{i'+L}
 \frac{\Delta(\widehat s_{i'})\Delta(v')}
 {\prod_{j\ne i'}\prod_{b<L}(s_j+v_b)}=(-1)^{i'+L}M_{i'}.
\]
So the residues agree, and the difference of the two sides is a polynomial
in $v_L$ of degree at most $h-1$, zero if $h=0$. It vanishes at the
$L-1\ge h$ points $v_L=v_b$, $b<L$, where the determinant has two equal
columns and $\Delta(v)=0$. Hence it is zero.

The product of the four factors is
$(-1)^{e_\Delta}a^{m(\sigma+k)+2e_\Delta}\Phi$. Since
$m=4n$, $e_\Delta=2n(4n-1)$ is even and
$m(\sigma+k)+2e_\Delta=3m^2/2-m$. All integrals converge absolutely:
$3-xv\ge2$ on $[0,1]^2$, $|z|\ge r$ on $\Gamma$, and $\Gamma$ has finite
variation.
\end{proof}

An exact calculation gives
\begin{equation}\label{eq:g4-polynomial}
G_4(T)=\frac{9}{67228000000000}\sum_{j=0}^4q_jT^j,\qquad
67228000000000=2^{11}5^97^5,
\end{equation}
with the coefficients in Table~\ref{tab:g4-coefficients}.
\begin{table}[!ht]
\centering
\caption{The integer coefficients in \eqref{eq:g4-polynomial}.}
\label{tab:g4-coefficients}
\begin{tabular}{cr}
$j$&$q_j$\\\hline
$0$ & $11390586457461030954563$\\
$1$ & $-124400937938521703661570$\\
$2$ & $509486272807705849983600$\\
$3$ & $-927382286191365446400000$\\
$4$ & $633018449291268096000000$\\
\end{tabular}
\end{table}

\section{The divisor}\label{sec:divisor}
Put
\[
N_{\mathrm{den}}=10n-1,\qquad U=11n-1,\qquad d_\ell=\operatorname{lcm}(1,\ldots,\ell),
 \qquad \mathcal D_m=d_{N_{\mathrm{den}}}^{3m}d_U^n.
\]
For $F\in\mathbb Q[T]$ let $v_p(F)=\min_j v_p([T^j]F)$. From now on $k$ is
a running index, and the parameter $k$ of Section~\ref{sec:determinant} equals
$n$. Fix the seam $R=\Rseam$. Primes $p\le m/R$ are cleared by
$\mathcal D_m$ alone. For primes $m/R<p<5m/4$ we use the rank bound of
Lemma~\ref{lem:residue-bound}.

\begin{lemma}\label{lem:common-denominator}
For an indeterminate $a$,
\begin{equation}\label{eq:lcm-clearer}
 P_m(a,T):=\mathcal D_mG_m(T;a)\in\Z[a,T].
\end{equation}
\end{lemma}
\begin{proof}
The denominators of a $C$ entry divide $d_{N_{\mathrm{den}}}$. An off-diagonal $B$ entry is
a difference of harmonic sums divided by $j-i$, so its denominator divides
$d_{N_{\mathrm{den}}}^2$, and so does that of $b_i$. The entries $Ta^i$ are
integral. The denominators of a $V$ entry divide $d_U$.
\end{proof}

\subsection{Residue classes}
For $\rho>0$ and $0\le x<\rho$ let
\begin{align}\label{eq:residue-sets}
 K(x,\rho)&=\{k\ge0:1/2\le x+k\rho<3/2\},\\
 \mathcal L(x,\rho)&=\{l\ge0:1/4\le x+l\rho<5/2\},\notag
\\
 W(x,\rho)&=\{w\ge1:0<w\rho-x\le1/4\},\notag
\end{align}
and
\begin{equation}\label{eq:residue-rank}
 \mathfrak r(x,\rho)=\min(2|K|+|W|,|\mathcal L|)
       -\mathbf1_{\{K=\{0\},\,\mathcal L=\{0,1\},\,W=\varnothing\}},
 \qquad \mathfrak c(\rho)=1+\int_0^\rho \mathfrak r(x,\rho)\,dx.
\end{equation}

\begin{lemma}\label{lem:residue-bound}
Let $p$ be a prime with $p^2>U$, so that $p\ge5$, and let $a=3$. Then
\begin{equation}\label{eq:inner-valuation}
 v_p(G_m)\ge-m-\sum_{u=0}^{p-1}\mathfrak r(u/m,p/m)=-m\mathfrak c(p/m).
\end{equation}
\end{lemma}
\begin{proof}
Multiply the whole rows of $C$, $B-TA$ and $V$ by $p$, $p^2$ and $p$. The
resulting matrix $S_p(T)$ has $\det S_p(T)=p^{3m+n}G_m(T)$ and entries in
$\Z_{(p)}[T]$, and $S_p(T)\equiv S_p(0)\pmod p$ because $p^2TA\equiv0$.
Write $i=u+kp$, $j=u+lp$ with
$0\le u<p$. Only the harmonic terms with $p$ in the denominator survive
modulo $p$. Their quotients are $p$-units because all indices are less than
$p^2$, and $a^p\equiv a\pmod p$, so
\[
 pt_{u+kp}\equiv a^ut_k,\qquad p^2b_{u+kp}\equiv a^ub_k\pmod p.
\]
The scaled $C$ and $B$ entries vanish modulo $p$ between distinct residue
classes modulo $p$. For a cross-class $B$ entry, the numerator is $p$-integral after
the row scaling and the remaining factor $p/(j-i)$ is zero modulo $p$.
Within a class, for $j\ne i$, $j-i=p(l-k)$ with $0<|l-k|<p$, so the
quotient denominator is a unit. Diagonal entries are given by the separate
cases $i=j$ in the definitions of $C$ and $B$. Since $a^{u+kp}\equiv a^{u+k}\pmod p$, for $i\ne j$
\[
 pC_{ij}=\frac{a^i-a^j}{l-k}\equiv a^u\frac{a^k-a^l}{l-k}=a^uC_{kl},\qquad
 p^2B_{ij}=\frac{pt_i-pt_j}{l-k}\equiv a^uB_{kl}\pmod p.
\]
A $V$ row reduces to zero unless $\ell+u=pw$. Then its
entry is $1/(l+w)$, without the factor $a^u$, and $l+w=(j+\ell)/p<p$. So,
after permutations, the reduction has blocks
with $2|K|+|W|$ rows and $|\mathcal L|$ columns, where $(x,\rho)=(u/m,p/m)$.

For $K=\{0\}$, $\mathcal L=\{0,1\}$, $W=\varnothing$ the block, with its rows
divided by the unit $a^u$, is
\begin{equation}\label{eq:rank-one-block}
 \begin{pmatrix}0&1-a\\0&-1\end{pmatrix}
 =\begin{pmatrix}0&-2\\0&-1\end{pmatrix},
\end{equation}
of rank one for every $p$. Every other block has rank at most the
minimum of its dimensions. So the rank of $S_p(0)$ modulo $p$ is at most
$\mathrm{rk}_p=\sum_u \mathfrak r(u/m,p/m)$.
The matrix has $2m+n$ rows. Row operations over $\Z_{(p)}$, independent of
$T$, make its last $2m+n-\mathrm{rk}_p$ rows divisible by $p$ coefficientwise.
Subtracting the scaling exponent $3m+n$ gives the first inequality in
\eqref{eq:inner-valuation}.

Every boundary in \eqref{eq:residue-sets} has the form $x=c_0+b\rho$ with
$4c_0,b\in\Z$. At $\rho=p/m$, $m=4n$, it lies in $m^{-1}\Z$. The $K,\mathcal L$
conditions are closed on the left and open on the right in $x$, and the
$W$ condition is $w\rho-1/4\le x<w\rho$. So $\mathfrak r$ is constant on
each $[u/m,(u+1)/m)$, and
\begin{equation}\label{eq:lattice-sum}
 \sum_{u=0}^{p-1}\mathfrak r(u/m,p/m)=m\int_0^{p/m}\mathfrak r(x,p/m)\,dx.\qedhere
\end{equation}
\end{proof}

The lines that bound the sets \eqref{eq:residue-sets} are
$x=z_0-k\rho$ $(z_0\in\{1/4,1/2,3/2,5/2\},\ k\ge0)$ and $x=w\rho-1/4$
$(w\ge1)$. The first meets the strip $0<x<\rho$ exactly for
$z_0/(k+1)<\rho<z_0/k$, the second exactly for $1/(4w)<\rho<1/(4(w-1))$. For
$k=0$ and for $w=1$ the upper endpoint is $+\infty$. So $x=z_0-k\rho$ meets
the strip for some $\rho\in(1/R,5/4)$ if and only if $4z_0/5-1<k<z_0R$, and
$x=w\rho-1/4$ if and only if $1\le w<R/4+1$. Hence, for $1/R<\rho<5/4$, the
sets \eqref{eq:residue-sets} change only on the $\nlines$ lines
\begin{equation}\label{eq:subdivision-lines}
\begin{gathered}
 x=0,\quad x=\rho,\quad x=w\rho-1/4\quad(1\le w\le\wmax),\\
 x=1/4-k\rho\quad(0\le k\le\kLlo),\quad x=1/2-k\rho\quad(0\le k\le\kKlo),\\
 x=3/2-k\rho\quad(1\le k\le\kKhi),\quad x=5/2-k\rho\quad(2\le k\le\kLhi).
\end{gathered}
\end{equation}
Cut $(1/R,5/4)$ at the finitely many values of $\rho$ where two of these
lines meet in $0\le x\le\rho$, and let $\mathcal C$ be the set of these cut
points. For $\rho\in(1/R,5/4)\setminus\mathcal C$ and each $z_0$ exactly one
line $x=z_0-k\rho$ meets $0<x<\rho$, namely $k=\lfloor z_0/\rho\rfloor$, and
exactly one line $x=w\rho-1/4$ does, namely $w=\lfloor1/(4\rho)\rfloor+1$.
These indices lie in the ranges of \eqref{eq:subdivision-lines}. So on each
open interval of $(1/R,5/4)\setminus\mathcal C$ the order of the lines in
$0<x<\rho$ is fixed, exactly five lines lie there, the sets
\eqref{eq:residue-sets} are constant on each of the six $x$-intervals
between consecutive lines, and $\mathfrak c$ is affine.
A change of $\rho$ by $\delta$ moves each of these lines by $O(|\delta|)$ and
the endpoint of the integral by $\delta$. So the sets
\eqref{eq:residue-sets} change only on strips of total length $O(|\delta|)$.
As $\mathfrak r$ is bounded,
$|\mathfrak c(\rho)-\mathfrak c(\rho')|=O(|\rho-\rho'|)$, and
$\mathfrak c$ is continuous. Table~\ref{tab:slabs}
(Appendix~\ref{app:arithmetic}) lists $\nslabs$ consecutive intervals with
union $(1/R,5/4)$ up to finitely many points, and the affine function
$\mathfrak c$ on each. Each row follows from
\eqref{eq:residue-sets}--\eqref{eq:residue-rank} and
\eqref{eq:subdivision-lines}. The affine functions of adjacent rows agree at
all $75$ internal joins. Summation of the rows gives
\begin{equation}\label{eq:inner-integral}
 \int_{1/R}^{5/4}\mathfrak c(\rho)\,d\rho=\innervalue.
\end{equation}
The configuration $K=\{0\}$, $\mathcal L=\{0,1\}$, $W=\varnothing$ of
\eqref{eq:rank-one-block} means $\max(1/2,3/2-\rho)\le x<3/2$,
$\max(1/4,5/2-2\rho)\le x<5/2-\rho$ and $x<\rho-1/4$. For $\rho<5/4$ this is
$\max(1/2,3/2-\rho,5/2-2\rho)\le x<\rho-1/4$, which is nonempty only for
$\rho>11/12$. So it occupies
$5/2-2\rho\le x<\rho-1/4$ for $11/12<\rho<1$, and
$1/2\le x<\rho-1/4$ for $1\le\rho<5/4$. As $1/R<11/12$, its contribution to
\eqref{eq:inner-integral} is
\begin{equation*}
 -\int_{11/12}^{1}(3\rho-11/4)\,d\rho
 -\int_1^{5/4}(\rho-3/4)\,d\rho=-\frac5{48}.
\end{equation*}

\subsection{Larger primes}
For $p>5n$ define row potentials $\lambda_C=0$, $\lambda_B(i)$,
$\lambda_V(\ell)$ and column potentials $\chi(j)$ by
\begin{equation*}
\begin{aligned}
5n<p<6n:&\quad
 (\lambda_B,\lambda_V,\chi)
 =(1+\mathbf1_{i\ge p},\,1,\,\mathbf1_{j\ge p+2n}),\\
6n\le p<8n:&\quad
 (\lambda_B,\lambda_V,\chi)
 =(\mathbf1_{i+p<10n},\,1,\,\mathbf1_{j\ge p}),\\
8n\le p<10n:&\quad
 (\lambda_B,\lambda_V,\chi)=(0,\,1,\,\mathbf1_{j\ge p}),\\
10n\le p\le11n-1:&\quad
 (\lambda_B,\lambda_V,\chi)=(0,\,\mathbf1_{p-\ell\in J},\,0),
\end{aligned}
\end{equation*}
and put
\begin{equation*}
 E_p^{\rm out}=
 \begin{cases}
 19n-2p,&5n<p<8n,\\
 11n-p,&8n\le p\le U.
 \end{cases}
\end{equation*}

\begin{lemma}\label{lem:outer-bound}
For a prime $5n<p\le U$ one has $v_p(G_m)\ge-E_p^{\rm out}$.
\end{lemma}
\begin{proof}
A $C$ entry has a pole only at $j=i+p$, of order at most one. A $B$ entry
has pole order at most
\[
 \mathbf1_{\max(i,j)\ge p}
 +\mathbf1_{\{\max(i,j)\ge p,\ i\equiv j\pmod p\}},
\]
including the double pole of $b_i$ for $i\ge p$. A $V$ entry has a simple
pole when $p\mid j+\ell$. As $p^2>U$, there are no higher powers.

In the first range every $B$ row has potential at least one, and two if
$i\ge p$. A collision $j=i+p$ has $j\ge p+2n$, and its column gives the
second unit. In the second range all $i<p$. The columns $j\ge p$ cover the
simple $B$ poles and the $C$ poles. A collision needs $i+p<10n$, and its
row gives the other unit. In the third range $i<p$ and $i+p\ge10n$, so
only simple $B$ poles remain. In these three ranges every $V$ row has
potential one, which covers $j+\ell=p$ and $j+\ell=2p$. In the fourth range
$i,j<p$, only $V$ poles remain, and the rows with $\lambda_V=1$ are those
containing one. So $\lambda_{\rm row}+\chi_j$ dominates every pole order.
Multiply each row and column by $p$ to the power of its potential. The
potentials add up to $(10n-p)+n+(8n-p)$, $(8n-p)+n+(10n-p)$, $n+(10n-p)$
and $11n-p$ in the four ranges.
\end{proof}

\subsection{The factor at three}
\begin{lemma}\label{lem:formal-divisor}
Let $N_m=\lceil53m^2/48\rceil-m$. Then
\begin{equation}\label{eq:formal-divisor}
 P_m(a,T)\in a^{N_m}\Z[a,T].
\end{equation}
In particular $v_3(G_m)\ge N_m-v_3(\mathcal D_m)$.
\end{lemma}
\begin{proof}
Work in $\mathbb Q[[a]][T]$. Write
\[
 S_j(a)=\sum_{u\ge1}\frac{a^u}{u^j},\qquad
 \mathsf v_k(i)=\begin{cases}(i-k)^{-1},&i\ne k,\\0,&i=k,\end{cases}
 \quad \mathsf z_k(i)=\mathsf v_k(i)^2,\quad \mathsf e_k(i)=\mathbf1_{i=k}.
\]
Subtract $S_1(a)$ times each $C$ row from the paired row of $B-TA$.
Comparison of coefficients, including the diagonal, gives
\begin{equation}\label{eq:formal-expansion}
\begin{aligned}
\begin{pmatrix}C\\B-TA-S_1C\end{pmatrix}
&=\sum_{k\ge0}a^k
 \begin{pmatrix}\mathsf v_k(I)\\-\mathsf z_k(I)-(T+S_2)\mathsf e_k(I)\end{pmatrix}\mathsf e_k(J)^t\\
&\quad+\sum_{k\ge0}a^k\begin{pmatrix}\mathsf e_k(I)\\\mathsf v_k(I)\end{pmatrix}\mathsf v_k(J)^t
 +\sum_{k\ge0}a^k\begin{pmatrix}0\\-\mathsf e_k(I)\end{pmatrix}\mathsf z_k(J)^t.
\end{aligned}
\end{equation}
For example, the off-diagonal lower entry is
$\sum_{k\ge0}a^k\mathsf v_k(i)\mathsf v_k(j)-(a^i+a^j)/(i-j)^2$, and the diagonal entry is
$\sum_{k\ge0}a^k\mathsf v_k(i)^2-a^i(T+S_2)=b_i-Ta^i$. Each coefficient of $a$
involves finitely many terms.

Expand along the $n$ rows of $V$, and expand each remaining $2m$-row minor
by the rank-one summands in \eqref{eq:formal-expansion}. A nonzero term
selects $2m=8n$ distinct rank-one summands. At weight $k$ at most
\begin{equation*}
 c_k=1+\mathbf1_{k\in J}+\mathbf1_{k\in I}
\end{equation*}
occur: the first summand needs $k\in J$, the third needs $k\in I$. All
entries, including $T+S_2$, have nonnegative order in $a$. So every term
has order at least the sum of the $8n$ smallest available weights.
The infinite expansion is used coefficientwise. Modulo
$a^{N_m+1}$ only finitely many $k$ occur, so Cauchy--Binet applies to a
finite truncation.

There is one weight at each $0\le k<n$, two at each $n\le k<2n$, and three
at each $2n\le k<6n$. Write $5n=3d+\epsilon'$, $0\le\epsilon'<3$. The minimum available-weight sum is
\begin{align*}
 \sum_{k=0}^{n-1}k+2\sum_{k=n}^{2n-1}k
 +3\sum_{k=2n}^{2n+d-1}k+\epsilon'(2n+d)
 &=\frac{53n^2}{3}-4n+\frac{\epsilon'(3-\epsilon')}6\\
 &=\left\lceil\frac{53n^2}{3}\right\rceil-4n=N_m.
\end{align*}
The minors of $V$ do not depend on $a,T$, so $a^{N_m}\mid G_m(T;a)$ in
$\mathbb Q[a,T]$. By \eqref{eq:lcm-clearer}, division by the monic
$a^{N_m}$ gives \eqref{eq:formal-divisor}. Specialization at $a=3$ gives
the last statement.
\end{proof}

\subsection{Summation over primes}
Put $e_0(p)=3m\lfloor\log_pN_{\mathrm{den}}\rfloor+n\lfloor\log_pU\rfloor=v_p(\mathcal D_m)$.
With the first applicable line taking precedence, let
\begin{equation}\label{eq:prime-exponents}
 E_p=\begin{cases}
 e_0(3)-N_m,&p=3,\\
 e_0(p),&p\le m/R\ \text{or}\ p^2\le U\ \text{or}\ p=5m/4,\\
 m\mathfrak c(p/m),&m/R<p<5m/4,\\
 E_p^{\rm out},&5m/4<p\le U,\\
 0,&p>U,
 \end{cases}
 \qquad
 K_m=\prod_{p}p^{E_p}\in\mathbb Q_{>0}.
\end{equation}

\begin{lemma}\label{lem:integerization}
One has $K_mG_m(T)\in\Z[T]$ for every $m=4n$, and
$\log K_m\le\tau m^2+o(m^2)$ as $m\to\infty$, where
\begin{equation}\label{eq:tau}
\begin{gathered}
 A=\frac{13}{4R}+\int_{1/R}^{5/4}\mathfrak c(\rho)\,d\rho+\frac{45}{32}
  =\Avalue,\\
 \tau=A-\frac{53}{48}\log3=\tauvalue\ldots,
\end{gathered}
\end{equation}
so that $\alpha_0-\tau>\marginvalue$ for the constant $\alpha_0$ defined after
\eqref{eq:alpha} in Section~\ref{sec:analytic}.
\end{lemma}
\begin{proof}
By \eqref{eq:lattice-sum}, $m\mathfrak c(p/m)\in\Z$. Lemmas~\ref{lem:common-denominator},
\ref{lem:residue-bound}, \ref{lem:outer-bound} and~\ref{lem:formal-divisor}
give $v_p(K_mG_m)\ge0$ for every $p$. The exponent at three may be
negative. Since $\mathfrak r\le|\mathcal L|$ and
$\sum_u|\mathcal L(u/m,p/m)|=|J|=9n$, every inner exponent is at most
$13n$.

If $m\ge11R^2/4$, then $p>m/R$ implies $p^2>m^2/R^2\ge11m/4>U$. Put
\[
 \eta_{\rm PNT}(X)=\sup_{y\ge X}|\theta_{\rm Ch}(y)/y-1|.
\]
By Fact~\ref{fact:pnt}, $\eta_{\rm PNT}(X)\to0$.
For an affine $f$ on $[t_0,t_1]\subset[1/R,11/4]$, Stieltjes integration by
parts gives
\begin{equation*}
 m\sum_{t_0m<p<t_1m} f(p/m)\log p
 =m^2\int_{t_0}^{t_1} f(\rho)\,d\rho
  +O_f\bigl(\eta_{\rm PNT}(m/R)m^2+m\log m\bigr).
\end{equation*}
There are finitely many pieces, and the primes at the finitely many points
of $\mathcal C\cup\{1/R,5/4,2,11/4\}$, at most one for each point and each
$m$, contribute $O(m\log m)$. For $p\le m/R$ the first powers contribute
$(13m/4)\theta_{\rm Ch}(m/R)=\frac{13}{4R}m^2+o(m^2)$, and the higher powers, including
those of two and three,
\[
 O\Bigl(m\sum_{j=2}^{\lfloor\log_2U\rfloor}\theta_{\rm Ch}(U^{1/j})\Bigr)
 =O(m^{3/2}\log(3m)).
\]
Subtract $N_m\log3$ once. For $p>5m/4$ the exponent is $m\mathfrak c_{\rm out}(p/m)$
with
\[
 \mathfrak c_{\rm out}(\rho)=\begin{cases}
 19/4-2\rho,&5/4<\rho<2,\\
 11/4-\rho,&2\le\rho\le11/4,
 \end{cases}
 \qquad
 \int_{5/4}^{11/4}\mathfrak c_{\rm out}(\rho)\,d\rho=\frac98+\frac9{32}=\frac{45}{32}.
\]
With \eqref{eq:inner-integral} this gives $A-\frac{53}{48}\log3=\tau$.
The total error is
$O(\eta_{\rm PNT}(m/R)m^2+m^{3/2}\log(3m))=o(m^2)$. The bounds
$\log3>2\sum_{k=0}^{99}((2k+1)2^{2k+1})^{-1}$ and
$\log3<2\sum_{k=0}^{99}((2k+1)2^{2k+1})^{-1}+8/(3\cdot201\cdot2^{201})$
give the digits of $\tau$ and of $\alpha_0-\tau$.
\end{proof}

\section{Analytic estimate}\label{sec:analytic}

In \eqref{eq:integral-identity} there are $m$, $L=5m/4$ and $m$ variables
$x$, $v$, $z$. Let
\[
 \mathcal K_x=\mathcal K_v=[0,1],\qquad
 \mathcal K_z=[1,3]\cup\{z:|z|=r\}.
\]
For finite positive measures put
\begin{align*}
 I(\sigma,\sigma')&=\iint\log|u-w|\,d\sigma(u)\,d\sigma'(w),\qquad
 D(\sigma,\sigma')=\iint\log|3-uw|\,d\sigma(u)\,d\sigma'(w),\\
 \mathcal Q(\mu)&=I(\mu_x,\mu_x)+I(\mu_v,\mu_v)+I(\mu_z,\mu_z)
                +I(\mu_v,\mu_z)-D(\mu_x,\mu_v).
\end{align*}
For a fixed $\nu=(\nu_x,\nu_v,\nu_z)$ let
$\mathcal U_i(q)=\int\log|q-t|\,d\nu_i(t)$,
$D_i(q)=\int\log|3-qt|\,d\nu_i(t)$ $(i=x,v)$,
\begin{align*}
 \psi_x&=2\mathcal U_x-D_v,\qquad
 \psi_v=2\mathcal U_v+\mathcal U_z-D_x,\qquad
 \psi_z=2\mathcal U_z+\mathcal U_v,\\
 F_x(q)&=\psi_x(q)+\tfrac34\log q,\quad
 F_v(q)=\psi_v(q)+\tfrac14\log q\quad(0<q\le1),\\
F_z(q)&=\psi_z(q)-\tfrac54\log|q|\quad(q\in\mathcal K_z).
\end{align*}
For a compact line segment or circle $\Gamma\subset\mathbb C$ let
$\ell_\Gamma$ be arclength on $\Gamma$. A segment meets the open disk
$B(q,t)=\{w:|w-q|<t\}$ in length at most $2t$. A circle meets it in
arclength at most $2\pi t$. If $t$
is at least the radius, this follows from the length of the circle. Otherwise the
intersection is an arc of at most a semicircle, of length at most $\pi/2$
times its chord. So
$\ell_\Gamma(B(q,t))\le2\pi t$ for all $q\in\mathbb C$, $t>0$.
A positive measure $\eta$ is \emph{tame} if
$\eta\le H(\ell_{\Gamma_1}+\dots+\ell_{\Gamma_N})$ for some $H\ge0$ and
compact segments or circles $\Gamma_1,\ldots,\Gamma_N$, that is, if $\eta$
is a finite sum of measures with arclength density at most $H$ on finitely
many segments and circles. Tame measures are finite and compactly
supported, and sums of tame measures are tame. Let $\mathfrak M$ be the set
of triples $\mu=(\mu_x,\mu_v,\mu_z)$ of tame measures of masses $(1,5/4,1)$
with $\operatorname{supp}\mu_x\subseteq\{|q|\le19/16\}$ and
$\operatorname{supp}\mu_v\subseteq\{\zeta\le|q|\le19/16\}$ for some
$\zeta>0$. On these supports $|3-uw|\ge3-(19/16)^2>1$.

\begin{lemma}\label{lem:concavity}
For $\mu,\nu\in\mathfrak M$, with $\psi_x,\psi_v,\psi_z$ formed from $\nu$,
\begin{equation}\label{eq:dual-inequality}
 \mathcal Q(\mu)\le-\mathcal Q(\nu)
            +\int\psi_x\,d\mu_x+\int\psi_v\,d\mu_v+\int\psi_z\,d\mu_z.
\end{equation}
\end{lemma}
\begin{proof}
We first show that for tame $\eta,\eta'$
\[
 \iint|\log|u-w||\,d\eta(u)\,d\eta'(w)<\infty.
\]
By the definition of tame it suffices to take $\eta=\ell_\Gamma$, $\eta'=\ell_{\Gamma'}$. Let
$\Lambda\ge1$ bound $|u-w|$ on $\Gamma\times\Gamma'$. For $u\in\Gamma$ the
layer-cake formula and the bound on $\ell_{\Gamma'}(B(u,t))$ give
\[
 \int|\log|u-w||\,d\ell_{\Gamma'}(w)
 \le\int_0^\infty\ell_{\Gamma'}(B(u,e^{-s}))\,ds+\ell_{\Gamma'}(\mathbb C)\log\Lambda
 \le2\pi+\ell_{\Gamma'}(\mathbb C)\log\Lambda,
\]
and integration over $u$ gives the bound.

Let $V$ be the space of real signed measures $\eta-\eta'$ with $\eta,\eta'$
tame and $\eta(\mathbb C)=\eta'(\mathbb C)$. By this bound $I$
is finite on $V\times V$ and on pairs of tame measures, so all pairings
below are finite. For $\sigma=\eta-\eta'\in V$ one has
$|\sigma|\le\eta+\eta'$, and Fact~\ref{fact:energy} gives
\[
 \|\sigma\|^2=-I(\sigma,\sigma)\ge0.
\]
So $-I$ is a positive semidefinite symmetric form on $V$, and
$|I(\sigma,\sigma')|\le\|\sigma\|\,\|\sigma'\|$.

Put $\sigma_i=\mu_i-\nu_i\in V$, $\iota(v)=3/v$ and $\beta=\iota_*\sigma_v$. The
$v$-parts are supported in $A=\{\zeta\le|q|\le19/16\}$ (the smaller
$\zeta$). The map $\iota$ sends $\Gamma\setminus\{0\}$ into a line or circle and
$A$ into $\{|q|\le3/\zeta\}$, and $|\iota'(v)|=3/|v|^2>1$ on $A$. So
$\iota_*(\ell_\Gamma|_A)\le\ell_{\Gamma'}$ for a compact segment or circle
$\Gamma'$. Hence $\iota_*\mu_v$ and $\iota_*\nu_v$ are tame, $\beta\in V$, and
$\log|v|$ is bounded on $A$. From
\begin{align*}
 \log|\iota(v)-\iota(w)|&=\log3+\log|v-w|-\log|v|-\log|w|,\\
 \log|3-xv|&=\log|x-\iota(v)|+\log|v|,
\end{align*}
and the zero masses,
\[
 I(\beta,\beta)=I(\sigma_v,\sigma_v),\qquad
 D(\sigma_x,\sigma_v)=I(\sigma_x,\beta).
\]
With $s_x=\|\sigma_x\|$, $s_v=\|\sigma_v\|=\|\beta\|$ and $s_z=\|\sigma_z\|$,
\begin{align*}
 \mathcal Q(\sigma)
 &=-s_x^2-s_v^2-s_z^2+I(\sigma_v,\sigma_z)-I(\sigma_x,\beta)
 \le-s_x^2-s_v^2-s_z^2+s_vs_z+s_xs_v\\
 &=-(s_x-s_v/2)^2-(s_z-s_v/2)^2-s_v^2/2\le0.
\end{align*}
Let $\mathcal Q(\cdot,\cdot)$ be the symmetric bilinear form with
$\mathcal Q(\mu,\mu)=\mathcal Q(\mu)$. Then
$2\mathcal Q(\nu,\mu)=\int\psi_x\,d\mu_x+\int\psi_v\,d\mu_v+\int\psi_z\,d\mu_z$
and
\[
 \mathcal Q(\mu)=\mathcal Q(\nu)+2\mathcal Q(\nu,\sigma)+\mathcal Q(\sigma)
 \le2\mathcal Q(\nu,\mu)-\mathcal Q(\nu).\qedhere
\]
\end{proof}

\begin{lemma}[computer-assisted]\label{lem:reference-bound}
The measure $\nu$ of Tables~\ref{tab:reference-intervals}
and~\ref{tab:reference-poisson} lies in $\mathfrak M$,
$\operatorname{supp}\nu_i\subseteq\mathcal K_i$, and
\begin{gather}
 \mathcal Q(\nu)\ge q_-:=\qminus,\label{eq:Q-lower}\\
 F_x\le K_x:=\Kx,\qquad F_v\le K_v:=\Kv,
 \qquad F_z\le K_z:=\Kz\label{eq:caps}
\end{gather}
on $(0,1]$, $(0,1]$ and $\mathcal K_z$. On $|z|=r$ even $F_z\le\Kzcirc$.
\end{lemma}
\begin{proof}
By Appendix~\ref{app:reference-data}, $\nu_x$, $\nu_v$ and the interval
part $\nu_z^{\rm int}$ of $\nu_z$ are sums of $131$, $143$ and $105$
intervals $[t_0,t_1]$ of positive mass $w$, each carrying
$w\,dt/(t_1-t_0)$. The intervals of each part form a chain of adjacent
intervals, contained in $[v_0,1]$ for $\nu_v$, $v_0=\vzero$. The circular
part $\nu_z^{\rm Poi}$ is a sum of $167$ Poisson measures $(c,w)$ with
$w>0$ and $|c|<r$,
\[
 \frac{w(r^2-|c|^2)}{2\pi|re^{i\vartheta}-c|^2}\,d\vartheta
 \qquad(0\le\vartheta\le2\pi).
\]
Its arclength density is at most $H_c$ in \eqref{eq:poisson-density}. The
masses are $1$, $5/4$ and $1$, and $\nu_z$ is invariant under
conjugation. Let $H_x,H_v,H_z$ be the largest values of $w/(t_1-t_0)$ over
the intervals of the three parts. So the parts of $\nu$ are tame,
$\operatorname{supp}\nu_v\subseteq[v_0,1]$, and $\nu\in\mathfrak M$. An
interval has potential
\[
 \frac{w}{t_1-t_0}\operatorname{Re}\{f(q-t_0)-f(q-t_1)\},\qquad f(u)=u\log u-u,
\]
with $u\log|u|-u$ on the real line and $0\log0=0$. A Poisson measure has potential
\begin{equation}\label{eq:poisson-potential}
 w\begin{cases}
 \log|r^2-q\overline c|-\log r,&|q|\le r,\\
 \log|q-c|,&|q|\ge r.
 \end{cases}
\end{equation}

\emph{Energy.} Interval pairs use $g(u)=\tfrac12u^2\log|u|-\tfrac34u^2$,
with $g''(u)=\log|u|$. A pair of Poisson measures $(c,w)$, $(c',w')$ gives
$ww'(\log|r^2-c\overline{c'}|-\log r)$, by \eqref{eq:poisson-potential} and
harmonicity. Interval--Poisson pairs integrate \eqref{eq:poisson-potential}
with the primitives of $\log|t-c|$ and $\log|r^2-t\overline c|$, split at
$t=r$. For the denominator, with $\mathsf M_j^i=\int t^j\,d\nu_i$ and
$N=\NE$,
\[
 D(\nu_x,\nu_v)=\frac54\log3-\sum_{j=1}^{N}\frac{\mathsf M_j^x\mathsf M_j^v}{j3^j}-\mathcal E_N,
 \qquad 0\le \mathcal E_N\le\frac{5/4}{2(N+1)3^{N}}.
\]
Table~\ref{tab:energy} gives the outward enclosures, at $\precE$-bit
precision. The enclosure of $D(\nu_x,\nu_v)$ contains the partial sum
minus $[0,(5/4)/(2(N+1)3^N)]$. The interval sum of the rows gives
$\mathcal Q(\nu)>q_-+5\cdot10^{-7}$.

\emph{Small arguments.} Put $\varepsilon_0=10^{-12}$. For $0<q\le\varepsilon_0$,
$\mathcal U_x(q),\mathcal U_v(q)\le0$, $\mathcal U_z(q)\le\log(3+\varepsilon_0)$,
$D_v(q)\ge\frac54\log(3-\varepsilon_0)$ and $D_x(q)\ge\log(3-\varepsilon_0)$. Hence
\begin{align*}
 F_x(q)&\le-\tfrac54\log(3-\varepsilon_0)+\tfrac34\log\varepsilon_0<-22.09,\\
 F_v(q)&\le\log(3+\varepsilon_0)-\log(3-\varepsilon_0)+\tfrac14\log\varepsilon_0<-6.90.
\end{align*}
Both are below the caps.

\emph{Subdivision.} The bounds on $[\varepsilon_0,1]$, $[\varepsilon_0,1]$,
$[1,3]$ and the upper half of $|z|=r$ are proved on the $\nleaves$ closed
pieces (leaves) of Lemma~\ref{lem:leaves}. The lower half of the circle
follows by conjugation. All bounds are computed in Arb ball arithmetic
\cite{Arb} at $\precF$-bit precision.

Let $[l,u]$ be a leaf on a real support, with centre $c_0$ and half-width
$h_0$. The interval components of $F$ come from the chains of $\nu_x$ on
$x$-leaves and of $\nu_v$, $\nu_z^{\rm int}$ on $v$- and $z$-leaves. The
endpoints of the chains of $\nu_x$, $\nu_v$, $\nu_z^{\rm int}$ cut the
initial pieces of the $x$-, $v$-, $z$-support. The chain of $\nu_z^{\rm int}$
misses the $v$-support, and that of $\nu_v$ meets the $z$-support only at
$1$. So the open leaf
meets at most one interval $[t_0,t_1]$ of each of these chains, and then lies
in it. A $v$-leaf lies in $[0,r]$ or in $[r,1]$.
Write $F=F_C+F_N$, where $F_C$ is the sum of the interval components of $F$,
with their coefficient $2$, whose interval contains $[l,u]$. Inside its
interval a component is convex, since $\mathcal U''=w/((q-t_0)(t_1-q))>0$
there, so $F_C$ lies below its chord on $[l,u]$. With
$B\ge\max\{0,\sup_{(l,u)}F_N''\}$, Taylor's formula with integral remainder
for $F_N$ at $c_0$ gives
\[
 \max_{[l,u]}F\le\max_{q\in\{l,u\}}\bigl(F_C(q)+F_N(c_0)+F_N'(c_0)(q-c_0)\bigr)
 +\tfrac12h_0^2B,
\]
because the chord plus an affine function is affine. The other interval
components have $\mathcal U''=-w/((q-t_0)(q-t_1))\le0$ on $(l,u)$. The terms
$\frac34\log q$ and $\frac14\log q$ are concave, and
$(-\frac54\log q)''\le\frac54l^{-2}$ on $[l,u]$. The terms $-D_v$, $-D_x$
are expanded as $-M\log3+\sum_{j\ge1}\mathsf M_jq^j/(j3^j)$ with the first
$N'=\NF$ moments. The omitted value is at most $M/(2(N'+1)3^{N'})$, the
omitted derivative lies in $[0,M/(2\cdot3^{N'})]$, and the second
derivative is at most $M/4$. On the real line a Poisson measure has second
derivative $w\operatorname{Re}(-\overline c^2/(r^2-q\overline c)^2)
\le w|c|^2/|r^2-q\overline c|^2$ for $|q|\le r$ and
$w\operatorname{Re}(-(q-c)^{-2})\le w/|q-c|^2$ for $|q|\ge r$; $B$ uses
enclosures of these second derivatives over the leaf.

On an arc $q=re^{i\vartheta}$, $|\vartheta-\vartheta_0|\le h_0$, of the
circle we take the smaller of two bounds. The first is
\[
 F(\vartheta_0)+h_0|F'(\vartheta_0)|+\tfrac12h_0^2B,\qquad
 B\ge\max\{0,\sup F''\},
\]
with derivatives in $\vartheta$ and
\[
 \frac{d}{d\vartheta}\log|q-s|=\operatorname{Re}\frac{iq}{q-s},\qquad
 \frac{d^2}{d\vartheta^2}\log|q-s|=\operatorname{Re}\frac{qs}{(q-s)^2}
 \le\frac{r|s|}{|q-s|^2}.
\]
Here $B$ uses enclosures over the arc of these second derivatives for the
Poisson terms ($s=c$). For an interval component of mass $w$ and density
$H$ on $[t_0,t_1]$ it uses the smaller upper end of two enclosures: of the
exact second derivative
$H\operatorname{Re}\bigl(q\log\frac{q-t_1}{q-t_0}+\frac{q^2}{q-t_1}-\frac{q^2}{q-t_0}\bigr)$,
and of $w\operatorname{Re}(qs/(q-s)^2)$
with $s$ ranging over $[t_0,t_1]$. The second bound uses that the distance from $re^{i\vartheta}$
to a point $s\ge0$ increases with $\vartheta\in[0,\pi]$. It bounds the
real-interval terms by their value at the right end of the arc and applies
the first bound to the Poisson terms alone. On the arc starting at
$\vartheta=0$ only the second bound is used.

On every leaf the upper end of the resulting bound is below the cap of its
support in Table~\ref{tab:caps}: $K_x$, $K_v$, $K_z$ on $[1,3]$, and
$\Kzcirc<K_z$ on the circle. Table~\ref{tab:caps} gives the largest upper
end on each support. The least margin, attained on $[1,3]$, is at least
$\capmargin$.
\end{proof}

For this $\nu$ put
\begin{equation}\label{eq:alpha}
 \alpha=q_--\tfrac32\log3-K_x-\tfrac54K_v-K_z.
\end{equation}
Since $\frac32\log3<1647919/10^6$,
\[
 \alpha>\alpha_0:=q_--\frac{1647919}{10^6}-K_x-\frac54K_v-K_z=\alphazero.
\]

\begin{lemma}\label{lem:potential-modulus}
Let $\eta$ be a compactly supported positive measure of mass $M$ with
$\eta(B(q,t))\le At$ for all $q\in\mathbb C$, $t>0$. If $0<d=|q-q'|\le1/3$, then
\[
 \Bigl|\int\log|q-t|\,d\eta(t)-\int\log|q'-t|\,d\eta(t)\Bigr|
 \le(6A+2M)\,d\,(1+\log(1/d)).
\]
The components of the reference of Lemma~\ref{lem:reference-bound} are
compactly supported, and for them one may take
\begin{equation}\label{eq:disk-constants}
 A_x=2H_x,\qquad A_v=2H_v,\qquad
 A_z=2H_z+\sum_c\frac{w_c(r+|c|)}{r(r-|c|)},
\end{equation}
the sum over the $167$ Poisson measures $(c,w_c)$. It is finite, since
$|c|<r$ for every centre.
\end{lemma}
\begin{proof}
The closed disk $\{w:|w-q|\le s\}$ is the intersection of the $B(q,s')$,
$s'>s$, so its $\eta$-measure is also at most $As$.
Let $N=\{t:|q-t|\le2d\}$. On $N$ both distances are at most $3d\le1$, and
the layer-cake formula gives
\begin{align*}
 \int_N\log\frac1{|q-t|}\,d\eta&\le2Ad\log\frac1{2d}+2Ad,\\
 \int_N\log\frac1{|q'-t|}\,d\eta&\le2Ad\log\frac1{3d}+3Ad.
\end{align*}
Off $N$, $|q'-t|\ge|q-t|/2$, the logarithmic difference is at most
$2d/|q-t|$, and
\[
 \int_{|q-t|>2d}\frac{d\eta(t)}{|q-t|}
 \le\int_0^1M\,ds+\int_1^{1/(2d)}\frac As\,ds
 =M+A\log\frac1{2d}.
\]
The sum is at most $6Ad\log(1/d)+5Ad+2Md$.

By the bounds on $\ell_\Gamma(B(q,t))$ before Lemma~\ref{lem:concavity}, an
interval of density at most $H_i$ gives $2H_i$, and a Poisson measure gives
$2\pi H_c$: since $|re^{i\vartheta}-c|\ge r-|c|$ and arclength is
$r\,d\vartheta$, its arclength density is at most
\begin{equation}\label{eq:poisson-density}
 H_c=\frac{w(r+|c|)}{2\pi r(r-|c|)}.\qedhere
\end{equation}
\end{proof}

Put $M_x=1$, $M_v=5/4$, $M_z=1$, and
\begin{equation}\label{eq:S-constant}
 B_i=6A_i+2M_i,\qquad
 S=2B_x+\tfrac72B_v+\tfrac{13}4B_z,\qquad C=10+3S.
\end{equation}
$S$ and $C$ are finite and depend only on $\nu$.

\begin{lemma}\label{lem:analytic-bound}
For every $m=4n\ge4$,
\begin{equation}\label{eq:analytic-full}
\begin{split}
 \log|G_m(\theta)|\le{}&-\alpha m^2+\tfrac{13}{2}m\log m
   +3S\bigl(1+\log(m^2/3)\bigr)+\tfrac{45}{4}\\
  &+m\log(T_r/3)-\log\bigl((m!)^2L!\bigr),
\end{split}
\end{equation}
and hence $\log|G_m(\theta)|\le-\alpha_0m^2+Cm\log m$ with $C$ as in
\eqref{eq:S-constant}.
\end{lemma}
\begin{proof}
Expand $\Gamma^m$ in \eqref{eq:integral-identity} into products of its
pieces and take absolute values. By \eqref{eq:chain-variation} the paths
contribute a factor at most $T_r^m$, and $[0,1]^m\times[0,1]^L$ has volume
one. So it suffices to bound $\log|\Phi|$ for $x\in[0,1]^m$,
$v\in[0,1]^L$, $z\in\mathcal K_z^m$. If $\Phi=0$ there is nothing to
prove. Otherwise $x_{i'},v_b>0$ and the points of each species are distinct.
Let $\mu_x$, $\mu_v$, $\mu_z$ put mass $1/m$ at each $x_{i'}$, $v_b$, $z_d$. Since
$\sigma+k+h=3m/4$, $\sigma=m/4$ and $m+k=5m/4$,
\begin{equation}\label{eq:energy-form}
 \frac{\log|\Phi|}{m^2}=\mathcal Q_{\rm disc}(\mu)
 +\frac34\int\log x\,d\mu_x+\frac14\int\log v\,d\mu_v
 -\frac54\int\log|z|\,d\mu_z,
\end{equation}
where $\mathcal Q_{\rm disc}$ omits the diagonal in the three self-energies.

Put $\epsilon=m^{-2}\le1/16$. Replace every atom $q$ by the uniform measure
of mass $1/m$ on the circle $q+2i\epsilon+\epsilon e^{i\vartheta}$, and let
$\mu^\epsilon$ be the result. Since
$\frac1{2\pi}\int_0^{2\pi}\log|u+\epsilon e^{i\vartheta}|\,d\vartheta
=\log\max(|u|,\epsilon)$ and the common translation cancels in
differences, averaging does not decrease any off-diagonal positive pair. Each
circle has energy $\log\epsilon$, and there are $13m/4$ atoms. The
function $\log|3-XV|$ is harmonic in each variable on the product of two
smearing disks, since $|X|,|V|\le1+3\epsilon\le19/16$. Its mean is
$\log|3-(x+2i\epsilon)(v+2i\epsilon)|\le\log(3-xv)+3\epsilon$, because
$|(x+2i\epsilon)(v+2i\epsilon)-xv|\le4\epsilon+4\epsilon^2$ and $3-xv\ge2$.
With the masses $1\cdot5/4$,
\begin{equation}\label{eq:smoothing}
 \mathcal Q_{\rm disc}(\mu)\le\mathcal Q(\mu^\epsilon)
 +\frac{13}{4m}\log\frac1\epsilon+\frac{15}4\epsilon.
\end{equation}
The measures $\mu^\epsilon$ and $\nu$ lie in $\mathfrak M$: $\mu^\epsilon$
is a sum of uniform measures on circles, inside $\{|q|\le19/16\}$ for $x$
and $v$, with $\operatorname{Im}q\ge\epsilon$ on the $v$-circles. For $\nu$
this is part of Lemma~\ref{lem:reference-bound}. By Lemmas~\ref{lem:concavity}
and~\ref{lem:reference-bound},
\[
 \mathcal Q(\mu^\epsilon)\le-q_-+\sum_{i}\int\psi_i\,d\mu_i^\epsilon.
\]

Each point of a smearing circle lies within $d=3\epsilon\le3/16$ of the
atom $q$ it replaces. Since $s\mapsto s(1+\log(1/s))$ increases on $(0,1]$,
Lemma~\ref{lem:potential-modulus} bounds the change of $\mathcal U_i$ by
$B_id(1+\log(1/d))$. For $q\in[0,1]$, $|p-q|\le d$ and
$0\le t\le1$, $|\nabla_p\log|3-pt||\le t/|3-pt|\le1$, so $D_i$ changes by
at most $M_id$. Counting masses, the coefficients of $B_x,B_v,B_z$ are
$2$, $\frac54\cdot2+1=\frac72$ and $\frac54+2=\frac{13}4$, and the $D$ terms
give $\frac54+\frac54=\frac52$. Hence
\[
 \sum_i\int\psi_i\,d\mu_i^\epsilon
 \le\sum_i\int\psi_i\,d\mu_i+Sd(1+\log(1/d))+\tfrac52d.
\]
The atoms lie in $(0,1]$, $(0,1]$ and $\mathcal K_z$. Adding the original
external terms of \eqref{eq:energy-form} turns $\psi_i$ into $F_i$, and
\eqref{eq:caps} gives
\[
 \frac{\log|\Phi|}{m^2}\le-q_-+K_x+\tfrac54K_v+K_z
 +\frac{13}{4m}\log\frac1\epsilon+3S\epsilon\Bigl(1+\log\frac1{3\epsilon}\Bigr)
 +\frac{45}4\epsilon,
\]
with $45/4=15/4+3\cdot5/2$. Multiply by $m^2$, put $\epsilon=m^{-2}$, add
the logarithm $(3m^2/2-m)\log3-\log((m!)^2L!)$ of the prefactor in
\eqref{eq:integral-identity} and $m\log T_r$. By \eqref{eq:alpha} this is
\eqref{eq:analytic-full}.

For the last statement drop $-\log((m!)^2L!)$ and $m\log(T_r/3)<0$. As
$m\ge4$, $3S(1+2\log m-\log3)\le6S\log m\le3Sm\log m$ and
$45/4\le3m\log m$. So the right-hand side is at most
$-\alpha m^2+(10+3S)m\log m=-\alpha m^2+Cm\log m$.
\end{proof}

\section{Nonvanishing}\label{sec:nonvanishing}

Let $\Z_{(17)}=\{u/v\in\mathbb Q:17\nmid v\}$, and let $v_{17}$ be the
$17$-adic valuation, $v_{17}(0)=+\infty$.

\begin{lemma}\label{lem:g-mod-17}
For $m=4$, coefficientwise in $\Z_{(17)}[T]$,
\begin{equation}\label{eq:g-mod-17}
G_4(T)\equiv g(T):=11+5T^2+4T^4\pmod {17}.
\end{equation}
\end{lemma}

\begin{proof}
Here $I=\{2,\ldots,5\}$, $J=\{1,\ldots,9\}$, $\Lambda=\{1\}$, and all
denominators divide $6350400=2520^2$, prime to $17$. Modulo $17$ the matrix
\eqref{eq:determinant} is
\begin{equation}\label{eq:g4-mod-matrix}
\setlength{\arraycolsep}{3.5pt}
\begin{pmatrix}
11&0&16&15&7&7&3&13&9\\
5&16&0&14&11&4&4&9&5\\
8&15&14&0&8&16&12&12&10\\
8&7&11&8&0&7&14&2&2\\
6&16+8T&4&13&10&3&15&14&8\\
5&4&16+7T&5&13&14&5&16&3\\
5&13&5&13+4T&4&10&5&6&6\\
9&10&13&4&3+12T&16&14&1&15\\
9&6&13&7&3&5&15&2&12
\end{pmatrix}.
\end{equation}
Its determinant is $g(T)$. Reduction of \eqref{eq:g4-polynomial}, where
$9/67228000000000\equiv13\pmod{17}$ and
$(q_0,q_1,q_2,q_3,q_4)\equiv(10,0,3,0,16)\pmod{17}$, gives the same
polynomial.
\end{proof}

\begin{lemma}\label{lem:nonvanishing}
Let $\xi\in\mathbb Q$. For every integer $t\ge1$ with $2t+v_{17}(\xi)>0$,
$G_{4\cdot17^t}(\xi)\ne0$.
\end{lemma}

\begin{proof}
Put $Q_t=17^t$, $m_t=4Q_t$ and $h=\sigma=k=Q_t$, and
\begin{equation}\label{eq:scaled-polynomial}
P_t(W)=Q_t^{3m_t+h}G_{m_t}(W/Q_t^2).
\end{equation}
This is the determinant obtained after $T=W/Q_t^2$ by multiplying the $C$,
$B-TA$ and $V$ rows by $Q_t$, $Q_t^2$ and $Q_t$. The middle rows become
$Q_t^2B-WA$. Every nonzero index,
difference and Cauchy denominator has absolute value less than $17Q_t$,
hence $17$-adic valuation at most $t$. So $Q_tt_i$ and $Q_t^2b_i$ lie in
$\Z_{(17)}$, and for $i\ne j$
\[
 Q_t^2B_{ij}=\frac{Q_t}{j-i}\,(Q_tt_i-Q_tt_j),\qquad
 Q_tC_{ij}=\frac{Q_t}{j-i}\,(a^i-a^j),\qquad
 Q_tV_{\ell j}=\frac{Q_t}{j+\ell},
\]
with $v_{17}(j-i),v_{17}(j+\ell)\le t$. So all entries of the scaled matrix
lie in $\Z_{(17)}[W]$. Hence
$P_t\in\Z_{(17)}[W]$.

Let $0\le u<Q_t$. Only the summation indices divisible by $Q_t$
survive, and $a^{Q_t}\equiv a\pmod{17}$, so
\begin{equation*}
Q_t t_{Q_t j+u}\equiv a^ut_j,
\qquad Q_t^2b_{Q_t j+u}\equiv a^ub_j\pmod {17}.
\end{equation*}
The scaled $C$ and $B$ entries vanish modulo $17$ between distinct residue
classes modulo $Q_t$. In the class $u$ write $i=Q_t i_0+u$, $j=Q_t j_0+u$. Then
$2\le i_0\le5$, $1\le j_0\le9$, and $A_{ii}\equiv a^uA_{i_0i_0}$. The only $V$ row in this class is
$\ell=Q_t-u$, with
\[
\frac{Q_t}{j+\ell}=\frac{1}{j_0+1}.
\]
Cross-class $C$ and $B$ entries have $17$-adic valuation at least $1$.
Within a class the nonzero differences $j_0-i_0$ satisfy $|j_0-i_0|\le7<17$
and are units. All other scaled $V$ entries vanish modulo $17$. So, after permutations,
the matrix has $Q_t$ diagonal blocks
\[
\begin{bmatrix}
a^u C_4\\ a^u(B_4-WA_4)\\ V_4
\end{bmatrix},\qquad 0\le u<Q_t,
\]
where the subscript $4$ refers to $G_4$. Each block has eight rows with the
factor $a^u$. Hence
\begin{equation}\label{eq:block-congruence}
P_t(W)\equiv\varepsilon_t\,3^{4Q_t(Q_t-1)}g(W)^{Q_t}\pmod {17},
\qquad \varepsilon_t\in\{1,-1\},
\end{equation}
coefficientwise. In particular $P_t(0)\equiv\pm3^{4Q_t(Q_t-1)}11^{Q_t}\not\equiv0\pmod{17}$.

Take $W=Q_t^2\xi$. Then $v_{17}(W)=2t+v_{17}(\xi)>0$, and since
$P_t\in\Z_{(17)}[W]$, $P_t(W)\equiv P_t(0)\not\equiv0\pmod{17}$. By
\eqref{eq:scaled-polynomial}, $G_{m_t}(\xi)\ne0$.
\end{proof}

The restriction $2t+v_{17}(\xi)>0$ is needed in this argument: $g$ has the
roots $7$ and $10$ in $\mathbb F_{17}$.

\section{Proof of Theorem~\ref{thm:main}}\label{sec:proof}

Suppose $\theta=u/q_0$ with $u\in\Z$ and $q_0\in\Z_{>0}$. By
Lemmas~\ref{lem:moment-identity} and~\ref{lem:integerization},
\[
Z_m:=q_0^mK_mG_m(u/q_0)\in\Z.
\]
By Lemma~\ref{lem:nonvanishing} with $\xi=\theta$, $Z_m\ne0$ for
$m=4\cdot17^t$ and all large $t$.
By Lemmas~\ref{lem:integerization} and~\ref{lem:analytic-bound},
\begin{equation*}
\log|Z_m|\le-(\alpha_0-\tau)m^2+Cm\log m+o(m^2)+m\log q_0.
\end{equation*}
As $\alpha_0-\tau>\marginvalue>0$, the right-hand side tends to $-\infty$ along this
sequence. So $0<|Z_m|<1$ for large $t$, which is impossible.
\qed

\appendix

\section{Arithmetic integration}\label{app:arithmetic}
Table~\ref{tab:line-families} lists the lines \eqref{eq:subdivision-lines}
for the seam $R=\Rseam$. The index ranges are the integer solutions of
$4z_0/5-1<k<z_0R$ and $1\le w<R/4+1$ (Section~\ref{sec:divisor}).

Table~\ref{tab:slabs} gives, for $\rho_0<\rho<\rho_1$, the function
$\mathfrak c(\rho)=c_0+c_1\rho$. Consecutive rows share an endpoint. The
first row starts at $1/R$ and the last ends at $5/4$. The rows are the
$\nslabs$ open intervals of $(1/R,5/4)\setminus\mathcal C$, so a row can be
checked at one interior point $\rho$: the six cells of
\eqref{eq:residue-sets} at $\rho$ give an affine function of $\rho$, which
must equal $c_0+c_1\rho$. Then
\[
 \int_{1/R}^{5/4}\mathfrak c(\rho)\,d\rho
 =\sum_{\rm rows}(\rho_1-\rho_0)\Bigl(c_0+c_1\frac{\rho_0+\rho_1}2\Bigr).
\]

\begin{center}\small\renewcommand{\arraystretch}{1.12}
\begin{minipage}{\textwidth}\centering
\captionof{table}{Line families of \eqref{eq:subdivision-lines}.}\label{tab:line-families}
\begin{tabular}{llr}
\hline
family & lines & count\\
\hline
strip & $x=0$, $x=\rho$ & $2$\\
$\mathcal L$ lower & $x=1/4-k\rho$, $0\le k\le\kLlo$ & $\the\numexpr\kLlo+1\relax$\\
$K$ lower & $x=1/2-k\rho$, $0\le k\le\kKlo$ & $\the\numexpr\kKlo+1\relax$\\
$K$ upper & $x=3/2-k\rho$, $1\le k\le\kKhi$ & $\kKhi$\\
$\mathcal L$ upper & $x=5/2-k\rho$, $2\le k\le\kLhi$ & $\the\numexpr\kLhi-1\relax$\\
$W$ & $x=w\rho-1/4$, $1\le w\le\wmax$ & $\wmax$\\
\hline
total & & $\nlines$\\
\hline
\end{tabular}
\end{minipage}
\end{center}

\begingroup\small
\setlength{\tabcolsep}{4pt}\renewcommand{\arraystretch}{1.12}
\begin{longtable}{@{}llll@{\hspace{1.4em}}llll@{\hspace{1.4em}}llll@{}}
\caption{The affine pieces of $\mathfrak c$ on $(1/R,5/4)$, $R=\Rseam$:
$\mathfrak c(\rho)=c_0+c_1\rho$ for $\rho_0<\rho<\rho_1$.}\label{tab:slabs}\\
\hline
$\rho_0$&$\rho_1$&$c_0$&$c_1$&$\rho_0$&$\rho_1$&$c_0$&$c_1$&$\rho_0$&$\rho_1$&$c_0$&$c_1$\\
\hline
\endfirsthead
\hline
$\rho_0$&$\rho_1$&$c_0$&$c_1$&$\rho_0$&$\rho_1$&$c_0$&$c_1$&$\rho_0$&$\rho_1$&$c_0$&$c_1$\\
\hline
\endhead
$1/8$ & $11/84$ & $9/2$ & $-10$ & $2/11$ & $11/60$ & $11/4$ & $2$ & $1/3$ & $11/32$ & $9/2$ & $-4$\\
$11/84$ & $5/38$ & $7/4$ & $11$ & $11/60$ & $3/16$ & $11/4$ & $2$ & $11/32$ & $7/20$ & $7/4$ & $4$\\
$5/38$ & $9/68$ & $17/4$ & $-8$ & $3/16$ & $5/26$ & $1/2$ & $14$ & $7/20$ & $5/14$ & $7/2$ & $-1$\\
$9/68$ & $2/15$ & $17/4$ & $-8$ & $5/26$ & $7/36$ & $3$ & $1$ & $5/14$ & $3/8$ & $6$ & $-8$\\
$2/15$ & $7/52$ & $9/4$ & $7$ & $7/36$ & $11/56$ & $19/4$ & $-8$ & $3/8$ & $11/28$ & $15/4$ & $-2$\\
$7/52$ & $3/22$ & $4$ & $-6$ & $11/56$ & $1/5$ & $2$ & $6$ & $11/28$ & $2/5$ & $15/4$ & $-2$\\
$3/22$ & $11/80$ & $5/2$ & $5$ & $1/5$ & $9/44$ & $3$ & $1$ & $2/5$ & $5/12$ & $7/4$ & $3$\\
$11/80$ & $5/36$ & $5/2$ & $5$ & $9/44$ & $5/24$ & $21/4$ & $-10$ & $5/12$ & $7/16$ & $1/2$ & $6$\\
$5/36$ & $9/64$ & $5/4$ & $14$ & $5/24$ & $11/52$ & $4$ & $-4$ & $7/16$ & $9/20$ & $9/4$ & $2$\\
$9/64$ & $1/7$ & $7/2$ & $-2$ & $11/52$ & $3/14$ & $4$ & $-4$ & $9/20$ & $11/24$ & $9/2$ & $-3$\\
$1/7$ & $11/76$ & $9/2$ & $-9$ & $3/14$ & $7/32$ & $5/2$ & $3$ & $11/24$ & $1/2$ & $7/4$ & $3$\\
$11/76$ & $7/48$ & $7/4$ & $10$ & $7/32$ & $2/9$ & $17/4$ & $-5$ & $1/2$ & $11/20$ & $19/4$ & $-3$\\
$7/48$ & $5/34$ & $7/2$ & $-2$ & $2/9$ & $9/40$ & $9/4$ & $4$ & $11/20$ & $9/16$ & $2$ & $2$\\
$5/34$ & $3/20$ & $6$ & $-19$ & $9/40$ & $5/22$ & $9/4$ & $4$ & $9/16$ & $7/12$ & $17/4$ & $-2$\\
$3/20$ & $11/72$ & $15/4$ & $-4$ & $5/22$ & $11/48$ & $19/4$ & $-7$ & $7/12$ & $5/8$ & $6$ & $-5$\\
$11/72$ & $2/13$ & $15/4$ & $-4$ & $11/48$ & $1/4$ & $2$ & $5$ & $5/8$ & $2/3$ & $19/4$ & $-3$\\
$2/13$ & $5/32$ & $7/4$ & $9$ & $1/4$ & $11/40$ & $9/2$ & $-5$ & $2/3$ & $11/16$ & $11/4$ & $0$\\
$5/32$ & $7/44$ & $1/2$ & $17$ & $11/40$ & $5/18$ & $7/4$ & $5$ & $11/16$ & $3/4$ & $11/4$ & $0$\\
$7/44$ & $9/56$ & $9/4$ & $6$ & $5/18$ & $9/32$ & $17/4$ & $-4$ & $3/4$ & $5/6$ & $1/2$ & $3$\\
$9/56$ & $11/68$ & $9/2$ & $-8$ & $9/32$ & $2/7$ & $17/4$ & $-4$ & $5/6$ & $7/8$ & $3$ & $0$\\
$11/68$ & $1/6$ & $7/4$ & $9$ & $2/7$ & $7/24$ & $9/4$ & $3$ & $7/8$ & $11/12$ & $19/4$ & $-2$\\
$1/6$ & $11/64$ & $19/4$ & $-9$ & $7/24$ & $3/10$ & $4$ & $-3$ & $11/12$ & $1$ & $19/4$ & $-2$\\
$11/64$ & $9/52$ & $2$ & $7$ & $3/10$ & $11/36$ & $5/2$ & $2$ & $1$ & $9/8$ & $15/4$ & $-1$\\
$9/52$ & $7/40$ & $17/4$ & $-6$ & $11/36$ & $5/16$ & $5/2$ & $2$ & $9/8$ & $5/4$ & $6$ & $-3$\\
$7/40$ & $5/28$ & $6$ & $-16$ & $5/16$ & $9/28$ & $5/4$ & $6$ & & & &\\
$5/28$ & $2/11$ & $19/4$ & $-9$ & $9/28$ & $1/3$ & $7/2$ & $-1$ & & & &\\
\hline
\end{longtable}
\endgroup

With \eqref{eq:inner-integral},
\[
 A=\frac{13}{\Rfour}+\int_{1/R}^{5/4}\mathfrak c(\rho)\,d\rho+\frac{45}{32},
\]
the constant of \eqref{eq:tau}. The coefficient of $\log3$ is $53/48$, and the
block \eqref{eq:rank-one-block} contributes $-5/48$ to
\eqref{eq:inner-integral}.

\section{Certified values}\label{app:numerical}

Table~\ref{tab:energy} gives the energy of $\nu$. Here $\nu_z^{\rm int}$ and
$\nu_z^{\rm Poi}$ are the interval and Poisson parts of $\nu_z$, $I_{\rm int}$ is the
energy of $\nu_z^{\rm int}$, $I_{\rm mix}=I(\nu_z^{\rm int},\nu_z^{\rm Poi})$, $I_{\rm Poi}$ is the
energy of $\nu_z^{\rm Poi}$, and
\[
 \mathcal Q(\nu)=I(\nu_x,\nu_x)+I(\nu_v,\nu_v)+I_{\rm int}+2I_{\rm mix}+I_{\rm Poi}
 +I(\nu_v,\nu_z^{\rm int})+I(\nu_v,\nu_z^{\rm Poi})-D(\nu_x,\nu_v).
\]
Every displayed interval is outward rounded, and the row of $D(\nu_x,\nu_v)$
contains the tail $\mathcal E_N$. The interval sum of the first eight rows,
with these coefficients, exceeds $q_-$, which is \eqref{eq:Q-lower}.
\begin{table}[!ht]
\centering\small\renewcommand{\arraystretch}{1.14}\setlength{\tabcolsep}{5pt}
\caption{The pair decomposition of $\mathcal Q(\nu)$ and its enclosures.}
\label{tab:energy}
\begin{tabular}{@{}lr@{}}
\hline term & outward enclosure\\\hline
$I(\nu_x,\nu_x)$ & \texttt{[-1.604087582665193,-1.604087582665192]}\\
$I(\nu_v,\nu_v)$ & \texttt{[-2.279694146532524,-2.279694146532523]}\\
$I_{\rm int}$ & \texttt{[-0.159655877431972,-0.159655877431971]}\\
$I_{\rm mix}$ & \texttt{[0.226965417324502,0.226965417324503]}\\
$I_{\rm Poi}$ & \texttt{[-0.495113146744583,-0.495113146744582]}\\
$I(\nu_v,\nu_z^{\rm int})$ & \texttt{[0.278231287125003,0.278231287125004]}\\
$I(\nu_v,\nu_z^{\rm Poi})$ & \texttt{[-0.191796788958401,-0.191796788958400]}\\
$D(\nu_x,\nu_v)$ & \texttt{[1.195386024949843,1.195386024949844]}\\
$\mathcal Q(\nu)$ & \texttt{[-5.193571445508505,-5.193571445508504]}\\
\hline
\end{tabular}
\end{table}

\begin{table}[!ht]
\centering\small\renewcommand{\arraystretch}{1.14}\setlength{\tabcolsep}{5pt}
\caption{The caps \eqref{eq:caps}: leaves of Lemma~\ref{lem:leaves}, the
largest upper bound over the leaves (rounded upward at $10^{-6}$) and the
cap. On $\mathcal K_z$ the cap is $K_z=\max(\Kz,\Kzcirc)=\Kz$.}
\label{tab:caps}
\begin{tabular}{@{}lrrr@{}}
\hline
support&leaves&upper bound&cap\\
\hline
$x\in[\varepsilon_0,1]$ & $67584$ & $-4.831488$ & $-4.83148$\\
$v\in[\varepsilon_0,1]$ & $37120$ & $-4.743190$ & $-4.74318$\\
$z\in[1,3]$ & $13568$ & $0.043178$ & $0.04318$\\
$|z|=r$, $\operatorname{Im}z\ge0$ & $8192$ & $0.042823$ & $0.04283$\\
\hline
\end{tabular}
\end{table}

\section{Reference measure and subdivision}\label{app:reference-data}
Terminating decimals are exact. Other values are printed as fractions $p/q$.

\subsection{Intervals}
Table~\ref{tab:reference-intervals} gives $\nu_x$, $\nu_v$ and
$\nu_z^{\rm int}$. Each is a chain of adjacent intervals. Its heading gives
the start point $b_0$, and its row $(k,b_k,w_k)$ is the interval
$[b_{k-1},b_k]$ carrying $w_k\,dt/(b_k-b_{k-1})$. The three chains have
$131$, $143$ and $105$ rows and end at $1$, $1$ and $3$. Outside its chain a
part vanishes. The masses are $1$, $5/4$ and $170340748021/500000000000$.
\begingroup\small
\setlength{\tabcolsep}{4pt}\renewcommand{\arraystretch}{1.12}
\begin{longtable}{@{}rll@{\hspace{1.2em}}ll@{\hspace{1.2em}}ll@{}}
\caption{Interval chains of the reference measure. A line gives the rows $(k,b_k,w_k)$, $(k+1,b_{k+1},w_{k+1})$ and $(k+2,b_{k+2},w_{k+2})$.}\label{tab:reference-intervals}\\
\hline
$k$&$b_k$&$w_k$&$b_{k+1}$&$w_{k+1}$&$b_{k+2}$&$w_{k+2}$\\
\hline
\endfirsthead
\hline
$k$&$b_k$&$w_k$&$b_{k+1}$&$w_{k+1}$&$b_{k+2}$&$w_{k+2}$\\
\hline
\endhead
\multicolumn{7}{@{}l}{$\nu_x$: $b_0=0.08426519384873$}\\
1 & $0.09379426$ & $0.004878668$ & $0.1033233$ & $0.005388618$ & $0.1137017$ & $0.006994053$\\
4 & $0.1240801$ & $0.007569004$ & $0.1352634$ & $0.008564016$ & $0.1464466$ & $0.008818097$\\
7 & $0.1583869$ & $0.009590098$ & $0.1703271$ & $0.009692372$ & $0.1829732$ & $0.0103232$\\
10 & $0.1956193$ & $0.01034485$ & $0.2089171$ & $0.01087642$ & $0.2222149$ & $0.0108577$\\
13 & $0.2361074$ & $0.01131139$ & $0.25$ & $0.01127106$ & $0.2644278$ & $0.01166085$\\
16 & $0.2788557$ & $0.01161365$ & $0.293757$ & $0.01194875$ & $0.3086583$ & $0.01190387$\\
19 & $0.3239693$ & $0.01219079$ & $0.3392803$ & $0.01215494$ & $0.3549354$ & $0.01239796$\\
22 & $0.3705905$ & $0.01237491$ & $0.3865227$ & $0.01257814$ & $0.4024548$ & $0.01257039$\\
25 & $0.4185959$ & $0.01273679$ & $0.4347369$ & $0.01274778$ & $0.4510177$ & $0.01287879$\\
28 & $0.4672984$ & $0.01290965$ & $0.4836492$ & $0.01300665$ & $0.5$ & $0.01305858$\\
31 & $0.5163508$ & $0.01312226$ & $0.5327016$ & $0.01319818$ & $0.5489823$ & $0.01322769$\\
34 & $0.5652631$ & $0.01332781$ & $0.5814041$ & $0.01332593$ & $0.5975452$ & $0.01345034$\\
37 & $0.6134773$ & $0.01341514$ & $0.6294095$ & $0.01356788$ & $0.6450646$ & $0.01349774$\\
40 & $0.6607197$ & $0.01367922$ & $0.6760307$ & $0.01357396$ & $0.6913417$ & $0.01378771$\\
43 & $0.706243$ & $0.01364148$ & $0.7211443$ & $0.01389263$ & $0.7355722$ & $0.01370571$\\
46 & $0.75$ & $0.01399346$ & $0.7638926$ & $0.01376403$ & $0.7777851$ & $0.01408874$\\
49 & $0.7910829$ & $0.01381373$ & $0.8043807$ & $0.01418846$ & $0.8170268$ & $0.01385703$\\
52 & $0.8296729$ & $0.01427975$ & $0.8416131$ & $0.01389719$ & $0.8535534$ & $0.01437142$\\
55 & $0.8647366$ & $0.01392421$ & $0.8759199$ & $0.0144671$ & $0.8862983$ & $0.01394331$\\
58 & $0.8966767$ & $0.0145547$ & $0.9062057$ & $0.01394895$ & $0.9157348$ & $0.0146563$\\
61 & $0.9243738$ & $0.01394465$ & $0.9330127$ & $0.01475476$ & $0.9407245$ & $0.01392136$\\
64 & $0.9484364$ & $0.01485533$ & $0.9551881$ & $0.01387691$ & $0.9619398$ & $0.01497938$\\
67 & $0.9677024$ & $0.01379006$ & $0.9734651$ & $0.01513013$ & $0.978214$ & $0.01364607$\\
70 & $0.9829629$ & $0.01534824607$ & $0.9864815$ & $0.01264469$ & $0.99$ & $0.01468005$\\
73 & $0.9901963$ & $0.0009304798$ & $0.9903926$ & $0.0008176946$ & $0.9923846$ & $0.009526088$\\
76 & $0.9943766$ & $0.01098753$ & $0.9950495$ & $0.004282778$ & $0.9957224$ & $0.004357794$\\
79 & $0.9962801$ & $0.003907897$ & $0.9968377$ & $0.004198966$ & $0.9975297$ & $0.005760908$\\
82 & $0.9982217$ & $0.006780532$ & $0.9985756$ & $0.003909506$ & $0.9989295$ & $0.004519428$\\
85 & $0.9989498$ & $0.0002978646$ & $0.9989702$ & $0.0002759106$ & $0.9989851$ & $0.0002057553$\\
88 & $0.999$ & $0.0002063948$ & $0.9990431$ & $0.0006106305$ & $0.9990862$ & $0.0006201$\\
91 & $0.9991731$ & $0.001304877$ & $0.9992599$ & $0.00137548$ & $0.9993488$ & $0.001492587$\\
94 & $0.9994377$ & $0.001607745$ & $0.9994512$ & $0.0002561717$ & $0.9994647$ & $0.0002462361$\\
97 & $0.9995672$ & $0.002058841$ & $0.9996696$ & $0.002355955$ & $0.9996767$ & $0.0001814995$\\
100 & $0.9996838$ & $0.0001633233$ & $0.999753$ & $0.00181967$ & $0.9998222$ & $0.002143887$\\
103 & $0.9998327$ & $0.0003687629$ & $0.9998432$ & $0.0003554893$ & $0.9998716$ & $0.001051055$\\
106 & $0.9999$ & $0.001186225$ & $0.9999219$ & $0.001022622$ & $0.9999438$ & $0.001204602$\\
109 & $0.9999515$ & $0.0004787153$ & $0.9999593$ & $0.0005156937$ & $0.9999638$ & $0.0003258994$\\
112 & $0.9999684$ & $0.000343669$ & $0.9999753$ & $0.000575319$ & $0.9999822$ & $0.000676004$\\
115 & $0.9999861$ & $0.0004312648$ & $0.99999$ & $0.0005066644$ & $0.9999922$ & $0.0003232857$\\
118 & $0.9999944$ & $0.0003798382$ & $0.9999956$ & $0.0002423703$ & $0.9999968$ & $0.0002847743$\\
121 & $0.9999975$ & $0.000181713$ & $0.9999982$ & $0.0002135055$ & $0.9999986$ & $0.0001362365$\\
124 & $0.999999$ & $0.0001600714$ & $0.9999992$ & $0.0001021401$ & $0.9999994$ & $0.0001200076$\\
127 & $0.9999996$ & $0.00007657429$ & $0.9999997$ & $0.00008996647$ & $0.9999998$ & $0.00005740427$\\
130 & $0.9999999$ & $0.0001104681$ & $1$ & $0.0003313966$ & &\\
\hline
\multicolumn{7}{@{}l}{$\nu_v$: $b_0=0.00960735979838$}\\
1 & $0.01332222$ & $0.001656734$ & $0.01703709$ & $0.00542487$ & $0.02178601$ & $0.008072268$\\
4 & $0.02653494$ & $0.008280399$ & $0.03229758$ & $0.009887851$ & $0.03806023$ & $0.009561701$\\
7 & $0.04481193$ & $0.01077175$ & $0.05156363$ & $0.01030386$ & $0.05927546$ & $0.01128366$\\
10 & $0.0669873$ & $0.01079812$ & $0.07562625$ & $0.01163129$ & $0.08426519$ & $0.01116664$\\
13 & $0.09379426$ & $0.01188784$ & $0.1033233$ & $0.01146128$ & $0.1137017$ & $0.01209265$\\
16 & $0.1240801$ & $0.01171582$ & $0.1352634$ & $0.01227961$ & $0.1464466$ & $0.0119444$\\
19 & $0.1583869$ & $0.01244886$ & $0.1703271$ & $0.01216045$ & $0.1829732$ & $0.01261611$\\
22 & $0.1956193$ & $0.01236727$ & $0.2089171$ & $0.0127843$ & $0.2222149$ & $0.0125745$\\
25 & $0.2361074$ & $0.01295137$ & $0.25$ & $0.01278215$ & $0.2644278$ & $0.01312539$\\
28 & $0.2788557$ & $0.01298995$ & $0.293757$ & $0.01330075$ & $0.3086583$ & $0.01319874$\\
31 & $0.3239693$ & $0.01347729$ & $0.3392803$ & $0.01340707$ & $0.3549354$ & $0.01365223$\\
34 & $0.3705905$ & $0.01360987$ & $0.3865227$ & $0.01381909$ & $0.4024548$ & $0.01378893$\\
37 & $0.4185959$ & $0.01398487$ & $0.4347369$ & $0.01401875$ & $0.4510177$ & $0.01417834$\\
40 & $0.4672984$ & $0.014227$ & $0.4836492$ & $0.01434849$ & $0.5$ & $0.01442008$\\
43 & $0.5163508$ & $0.01450478$ & $0.5327016$ & $0.01460288$ & $0.5489823$ & $0.01464888$\\
46 & $0.5652631$ & $0.01477283$ & $0.5814041$ & $0.01478305$ & $0.5975452$ & $0.01493237$\\
49 & $0.6134773$ & $0.01490385$ & $0.6294095$ & $0.01508317$ & $0.6450646$ & $0.01501368$\\
52 & $0.6607197$ & $0.01522317$ & $0.6760307$ & $0.01511253$ & $0.6913417$ & $0.01535607$\\
55 & $0.706243$ & $0.01519763$ & $0.7211443$ & $0.01548108$ & $0.7355722$ & $0.01527532$\\
58 & $0.75$ & $0.01559771$ & $0.7638926$ & $0.01534276$ & $0.7777851$ & $0.0157047$\\
61 & $0.7910829$ & $0.01539736$ & $0.8043807$ & $0.01581345$ & $0.8170268$ & $0.01544195$\\
64 & $0.8296729$ & $0.01591008$ & $0.8416131$ & $0.0154806$ & $0.8535534$ & $0.01600481$\\
67 & $0.8647366$ & $0.01550264$ & $0.8759199$ & $0.01610211$ & $0.8862983$ & $0.01551437$\\
70 & $0.8966767$ & $0.01618908$ & $0.9062057$ & $0.01551011$ & $0.9157348$ & $0.01629058$\\
73 & $0.9243738$ & $0.01549424$ & $0.9330127$ & $0.01638817$ & $0.9407245$ & $0.01545718$\\
76 & $0.9484364$ & $0.0164881$ & $0.9551881$ & $0.01539698$ & $0.9619398$ & $0.0166144$\\
79 & $0.9677024$ & $0.01529057$ & $0.9734651$ & $0.01677105$ & $0.978214$ & $0.01512197$\\
82 & $0.9829629$ & $0.01700337185$ & $0.9864815$ & $0.01400531$ & $0.99$ & $0.01625595$\\
85 & $0.9901963$ & $0.001030202$ & $0.9903926$ & $0.0009054735$ & $0.9923846$ & $0.01054699$\\
88 & $0.9943766$ & $0.0121635$ & $0.9950495$ & $0.004740679$ & $0.9957224$ & $0.004823687$\\
91 & $0.9962801$ & $0.004325506$ & $0.9968377$ & $0.004647527$ & $0.9975297$ & $0.006376076$\\
94 & $0.9982217$ & $0.007504222$ & $0.9985756$ & $0.004326653$ & $0.9989295$ & $0.005001531$\\
97 & $0.9989498$ & $0.0003296327$ & $0.9989702$ & $0.0003053404$ & $0.9989851$ & $0.0002277013$\\
100 & $0.999$ & $0.0002284092$ & $0.9990431$ & $0.000675759$ & $0.9990862$ & $0.0006862373$\\
103 & $0.9991731$ & $0.001444043$ & $0.9992599$ & $0.001522167$ & $0.9993488$ & $0.001651755$\\
106 & $0.9994377$ & $0.001779182$ & $0.9994512$ & $0.0002834868$ & $0.9994647$ & $0.0002724926$\\
109 & $0.9995672$ & $0.002278363$ & $0.9996696$ & $0.002607138$ & $0.9996767$ & $0.0002008494$\\
112 & $0.9996838$ & $0.0001807361$ & $0.999753$ & $0.002013666$ & $0.9998222$ & $0.002372435$\\
115 & $0.9998327$ & $0.0004080737$ & $0.9998432$ & $0.0003933854$ & $0.9998716$ & $0.001163099$\\
118 & $0.9999$ & $0.001312675$ & $0.9999219$ & $0.001131631$ & $0.9999438$ & $0.001333007$\\
121 & $0.9999515$ & $0.0005297437$ & $0.9999593$ & $0.0005706635$ & $0.9999638$ & $0.0003606381$\\
124 & $0.9999684$ & $0.0003803018$ & $0.9999753$ & $0.0006366439$ & $0.9999822$ & $0.0007480609$\\
127 & $0.9999861$ & $0.0004772341$ & $0.99999$ & $0.0005606707$ & $0.9999922$ & $0.0003577452$\\
130 & $0.9999944$ & $0.0004203256$ & $0.9999956$ & $0.000268205$ & $0.9999968$ & $0.0003151286$\\
133 & $0.9999975$ & $0.000201082$ & $0.9999982$ & $0.0002362633$ & $0.9999986$ & $0.0001507581$\\
136 & $0.999999$ & $0.0001771335$ & $0.9999992$ & $0.0001130273$ & $0.9999994$ & $0.0001327994$\\
139 & $0.9999996$ & $0.00008473636$ & $0.9999997$ & $0.00009955619$ & $0.9999998$ & $0.0000635229$\\
142 & $0.9999999$ & $0.000122243$ & $1$ & $0.0003667206$ & &\\
\hline
\multicolumn{7}{@{}l}{$\nu_z^{\rm int}$: $b_0=28303039933611/17500000000000$}\\
1 & $1.647939$ & $0.000800039$ & $1.678561$ & $0.001634018$ & $1.709871$ & $0.002195722$\\
4 & $1.741181$ & $0.002569887$ & $1.773045$ & $0.002941359$ & $1.80491$ & $0.003206511$\\
7 & $1.837192$ & $0.003492292$ & $1.869474$ & $0.003707447$ & $1.902035$ & $0.003935108$\\
10 & $1.934597$ & $0.004118914$ & $1.967298$ & $0.00430909$ & $2$ & $0.004469524$\\
13 & $2.032702$ & $0.00462588$ & $2.065403$ & $0.004777443$ & $2.097965$ & $0.004901502$\\
16 & $2.130526$ & $0.005045769$ & $2.162808$ & $0.005145023$ & $2.19509$ & $0.005284974$\\
19 & $2.226955$ & $0.005358049$ & $2.258819$ & $0.005500629$ & $2.290129$ & $0.005547422$\\
22 & $2.321439$ & $0.005693562$ & $2.352061$ & $0.005716104$ & $2.382683$ & $0.005869112$\\
25 & $2.412486$ & $0.0058643$ & $2.442289$ & $0.006028879$ & $2.471144$ & $0.00599805$\\
28 & $2.5$ & $0.00617374$ & $2.527785$ & $0.0061167$ & $2.55557$ & $0.006304648$\\
31 & $2.582166$ & $0.00622033$ & $2.608761$ & $0.00642816$ & $2.634054$ & $0.006311412$\\
34 & $2.659346$ & $0.006538253$ & $2.683226$ & $0.006392705$ & $2.707107$ & $0.006640591$\\
37 & $2.729473$ & $0.006459666$ & $2.75184$ & $0.006737854$ & $2.772597$ & $0.006515946$\\
40 & $2.793353$ & $0.006824422$ & $2.812411$ & $0.006559483$ & $2.83147$ & $0.006912101$\\
43 & $2.848748$ & $0.006592625$ & $2.866025$ & $0.006992728$ & $2.881449$ & $0.006611407$\\
46 & $2.896873$ & $0.007069417$ & $2.910376$ & $0.006615346$ & $2.92388$ & $0.007152991$\\
49 & $2.935405$ & $0.006594439$ & $2.94693$ & $0.00724478$ & $2.956428$ & $0.006543148$\\
52 & $2.965926$ & $0.007358282$ & $2.973356$ & $0.00643243$ & $2.980785$ & $0.007538103528$\\
55 & $2.985393$ & $0.00532701$ & $2.99$ & $0.006421216$ & $2.990722$ & $0.001161462$\\
58 & $2.991445$ & $0.001127666$ & $2.992911$ & $0.002498185$ & $2.994377$ & $0.002795714$\\
61 & $2.995607$ & $0.00261479$ & $2.996838$ & $0.003080948$ & $2.997348$ & $0.001436316$\\
64 & $2.997859$ & $0.001592602$ & $2.9979$ & $0.0001394852$ & $2.99794$ & $0.000132549$\\
67 & $2.998056$ & $0.0003909291$ & $2.998172$ & $0.0004039291$ & $2.998197$ & $0.00008733582$\\
70 & $2.998222$ & $0.00008626531$ & $2.998371$ & $0.0005458644$ & $2.99852$ & $0.0005686528$\\
73 & $2.998725$ & $0.0008328935$ & $2.998929$ & $0.0009109751$ & $2.998965$ & $0.0001665063$\\
76 & $2.999$ & $0.0001617129$ & $2.99917$ & $0.0008451466$ & $2.999339$ & $0.0009489754$\\
79 & $2.999388$ & $0.0002982447$ & $2.999438$ & $0.0003002494$ & $2.999561$ & $0.0008270791$\\
82 & $2.999684$ & $0.0009730542$ & $2.999685$ & $0.00001266974$ & $2.999686$ & $0.000009434384$\\
85 & $2.999754$ & $0.0006105242$ & $2.999822$ & $0.0007153116$ & $2.999861$ & $0.0004647103$\\
88 & $2.9999$ & $0.0005473779$ & $2.999909$ & $0.0001464355$ & $2.999919$ & $0.0001487465$\\
91 & $2.999931$ & $0.0002195401$ & $2.999944$ & $0.000241353$ & $2.999956$ & $0.0002613112$\\
94 & $2.999968$ & $0.0003071303$ & $2.999975$ & $0.0001959502$ & $2.999982$ & $0.0002302158$\\
97 & $2.999986$ & $0.0001468951$ & $2.99999$ & $0.0001725934$ & $2.999992$ & $0.0001101306$\\
100 & $2.999994$ & $0.0001293995$ & $2.999996$ & $0.00008256952$ & $2.999997$ & $0.00009701663$\\
103 & $2.999998$ & $0.00006190621$ & $2.999999$ & $0.0001191514$ & $3$ & $0.0003570495$\\
\hline
\end{longtable}
\endgroup

\subsection{Poisson measures}
Table~\ref{tab:reference-poisson} gives $\nu_z^{\rm Poi}$. A row $(c,w)$
with $\operatorname{Im}c>0$ stands for the two Poisson measures on
$|q|=r=2/5$ with centres $c$ and $\overline c$, each of mass $w$. The row
$k=0$, with $c$ real, stands for one Poisson measure. The measures are those
of the proof of Lemma~\ref{lem:reference-bound}, and their logarithmic
potential is \eqref{eq:poisson-potential}. This gives $2\cdot83+1=167$
measures of total mass $329659251979/500000000000$.
\begingroup\small
\setlength{\tabcolsep}{3.5pt}\renewcommand{\arraystretch}{1.12}
\begin{longtable}{@{}rlll@{\hspace{1.5em}}rlll@{}}
\caption{Poisson centres $c$ and masses $w$: rows $k\ge1$ have $\operatorname{Im}c>0$, in increasing order of $\arg c$, and stand also for $\overline c$; row $k=0$ is the real centre.}\label{tab:reference-poisson}\\
\hline
$k$&$\operatorname{Re}c$&$\operatorname{Im}c$&$w$&$k$&$\operatorname{Re}c$&$\operatorname{Im}c$&$w$\\
\hline
\endfirsthead
\hline
$k$&$\operatorname{Re}c$&$\operatorname{Im}c$&$w$&$k$&$\operatorname{Re}c$&$\operatorname{Im}c$&$w$\\
\hline
\endhead
0 & $-0.394$ & $0$ & $0.006687968$ & 42 & $-0.2046129$ & $0.341376$ & $0.0002611643$\\
1 & $0.2535036$ & $0.308895$ & $0.0002724889$ & 43 & $-0.2080587$ & $0.3345857$ & $0.004943139$\\
2 & $0.2370883$ & $0.3196766$ & $0.001734768001$ & 44 & $-0.2188947$ & $0.327599$ & $0.006091431004$\\
3 & $0.221117$ & $0.3309249$ & $0.003160258$ & 45 & $-0.2294962$ & $0.3202616$ & $0.004978119$\\
4 & $0.2046129$ & $0.341376$ & $0.002355698$ & 46 & $-0.2370883$ & $0.3196766$ & $0.0003172227$\\
5 & $0.197$ & $0.341214$ & $0.00203316$ & 47 & $-0.239852$ & $0.3125812$ & $0.004953756$\\
6 & $0.1876159$ & $0.3510047$ & $0.001945593$ & 48 & $-0.249951$ & $0.3045661$ & $0.006293054005$\\
7 & $0.1742617$ & $0.3533679$ & $0.00333003$ & 49 & $-0.2597823$ & $0.2962249$ & $0.00498701$\\
8 & $0.1626066$ & $0.3588803$ & $0.003486625$ & 50 & $-0.2672805$ & $0.2948985$ & $0.0003407953$\\
9 & $0.152308$ & $0.3677041$ & $0.002375737$ & 51 & $-0.2693354$ & $0.2875664$ & $0.005025951$\\
10 & $0.1387865$ & $0.3687469$ & $0.003817593$ & 52 & $-0.2786001$ & $0.2786001$ & $0.006379229002$\\
11 & $0.1266471$ & $0.3730905$ & $0.003942579$ & 53 & $-0.2875664$ & $0.2693354$ & $0.005074617$\\
12 & $0.1155333$ & $0.3808623$ & $0.002775625$ & 54 & $-0.2948985$ & $0.2672805$ & $0.0003451508$\\
13 & $0.1019747$ & $0.3805748$ & $0.004181110002$ & 55 & $-0.2962249$ & $0.2597823$ & $0.005068562$\\
14 & $0.08946805$ & $0.3837075$ & $0.004331115$ & 56 & $-0.3045661$ & $0.249951$ & $0.006456653$\\
15 & $0.07764595$ & $0.3903525$ & $0.003082157$ & 57 & $-0.3125812$ & $0.239852$ & $0.005132843$\\
16 & $0.06418082$ & $0.3887375$ & $0.004515421$ & 58 & $-0.3196766$ & $0.2370883$ & $0.0003275158$\\
17 & $0.05142732$ & $0.3906293$ & $0.004618404$ & 59 & $-0.3202616$ & $0.2294962$ & $0.00520169$\\
18 & $0.03901082$ & $0.3960835$ & $0.003346546$ & 60 & $-0.327599$ & $0.2188947$ & $0.006408025$\\
19 & $0.02576883$ & $0.3931564$ & $0.004623723$ & 61 & $-0.3345857$ & $0.2080587$ & $0.005265605$\\
20 & $0.01289132$ & $0.393789$ & $0.005370324$ & 62 & $-0.341376$ & $0.2046129$ & $0.0002748567$\\
21 & $0$ & $0.398$ & $0.002954771$ & 63 & $-0.341214$ & $0.197$ & $0.005351403$\\
22 & $-0.01289132$ & $0.393789$ & $0.00558512$ & 64 & $-0.347477$ & $0.1857303$ & $0.006322745001$\\
23 & $-0.02576883$ & $0.3931564$ & $0.004147672$ & 65 & $-0.3533679$ & $0.1742617$ & $0.00539497$\\
24 & $-0.03861875$ & $0.3921028$ & $0.005377501$ & 66 & $-0.3597877$ & $0.1701669$ & $0.0002230895$\\
25 & $-0.05142732$ & $0.3906293$ & $0.004364202$ & 67 & $-0.3588803$ & $0.1626066$ & $0.005494093$\\
26 & $-0.05839873$ & $0.3936923$ & $0.0002681214$ & 68 & $-0.3640085$ & $0.1507773$ & $0.006236832$\\
27 & $-0.06418082$ & $0.3887375$ & $0.004435699$ & 69 & $-0.3687469$ & $0.1387865$ & $0.005515612$\\
28 & $-0.07686559$ & $0.3864294$ & $0.00549873$ & 70 & $-0.3747345$ & $0.1340822$ & $0.0002252512$\\
29 & $-0.08946805$ & $0.3837075$ & $0.004569928$ & 71 & $-0.3730905$ & $0.1266471$ & $0.005431967$\\
30 & $-0.09670611$ & $0.3860724$ & $0.0002359134$ & 72 & $-0.3770345$ & $0.1143722$ & $0.006398725$\\
31 & $-0.1019747$ & $0.3805748$ & $0.004673523$ & 73 & $-0.3805748$ & $0.1019747$ & $0.005431584$\\
32 & $-0.1143722$ & $0.3770345$ & $0.005552904$ & 74 & $-0.3860724$ & $0.09670611$ & $0.0002813461$\\
33 & $-0.1266471$ & $0.3730905$ & $0.004783363$ & 75 & $-0.3837075$ & $0.08946805$ & $0.005361819$\\
34 & $-0.1340822$ & $0.3747345$ & $0.0001976397$ & 76 & $-0.3864294$ & $0.07686559$ & $0.00656059$\\
35 & $-0.1387865$ & $0.3687469$ & $0.004891999$ & 77 & $-0.3887375$ & $0.06418082$ & $0.005341411$\\
36 & $-0.1507773$ & $0.3640085$ & $0.005599664001$ & 78 & $-0.3936923$ & $0.05839873$ & $0.000338135$\\
37 & $-0.1626066$ & $0.3588803$ & $0.00497547$ & 79 & $-0.3906293$ & $0.05142732$ & $0.005290732$\\
38 & $-0.1701669$ & $0.3597877$ & $0.0002060286$ & 80 & $-0.3921028$ & $0.03861875$ & $0.00668558776$\\
39 & $-0.1742617$ & $0.3533679$ & $0.004926876$ & 81 & $-0.3931564$ & $0.02576883$ & $0.005271869$\\
40 & $-0.1857303$ & $0.347477$ & $0.005847720003$ & 82 & $-0.3975206$ & $0.01952893$ & $0.0003596048$\\
41 & $-0.197$ & $0.341214$ & $0.004984704$ & 83 & $-0.393789$ & $0.01289132$ & $0.005295008$\\
\hline
\end{longtable}
\endgroup

\subsection{Dyadic subdivision}\label{app:depths}
The four supports are
\[
 s=x,v:\ [\varepsilon_0,1],\qquad s=z:\ [1,3],\qquad
 s=\circ:\ \{re^{i\pi t}:0\le t\le1\},
\]
the last parametrized by $t\in[0,1]$. The \emph{initial pieces} of a support
are the closed intervals between consecutive points of the following sets,
in increasing order:
\begin{itemize}
\item[$x$:] $\varepsilon_0$, $1$, and the endpoints of the $\nu_x$-chain in
$(\varepsilon_0,1)$;
\item[$v$:] $\varepsilon_0$, $1$, $r=2/5$, and the endpoints of the
$\nu_v$-chain in $(\varepsilon_0,1)$;
\item[$z$:] $1$, $3$, and the endpoints of the $\nu_z^{\rm int}$-chain in $(1,3)$;
\item[$\circ$:] $t=k/\Narcs$, $0\le k\le\Narcs$.
\end{itemize}
Since the chains start at $0.0842\ldots$, $v_0$ and $1.617\ldots$ and end at
$1$, $1$ and $3$, and $r$ is not an endpoint of the $\nu_v$-chain, there are
$132$, $145$, $106$ and $\Narcs$ initial pieces.

\begin{lemma}\label{lem:leaves}
Let $d_x=9$, $d_v=8$, $d_z=7$, $d_\circ=8$. Bisect every initial piece of the
support $s$ to depth $d_s$: a piece $[A,B]$ gives the $2^{d_s}$ leaves
\[
 [A+(B-A)(j-1)2^{-d_s},\,A+(B-A)j2^{-d_s}],\qquad 1\le j\le2^{d_s},
\]
in $q$ on the real supports and in $t$ on the circle. The leaves of each
support have rational endpoints, disjoint interiors and union equal to the
support. There are $132\cdot2^9$, $145\cdot2^8$, $106\cdot2^7$ and
$\Narcs\cdot2^8$ leaves, $\nleaves$ in all (Table~\ref{tab:depths}).
\end{lemma}

\begin{proof}
Within a piece the leaves are consecutive and their lengths add up to
$B-A$. The pieces of a support are consecutive and cover it. The counts are
products.
\end{proof}

The leaves are the closed pieces of the proof of
Lemma~\ref{lem:reference-bound}. On each of them the bound of that proof is
below the cap, so no piece needs a finer subdivision.

\begin{table}[!ht]
\centering\small\renewcommand{\arraystretch}{1.14}
\caption{Uniform bisection depths.}\label{tab:depths}
\begin{tabular}{@{}lrrr@{}}
\hline
support&initial pieces&depth $d_s$&leaves\\
\hline
$x$ & $132$ & $9$ & $67584$\\
$v$ & $145$ & $8$ & $37120$\\
$z$ real & $106$ & $7$ & $13568$\\
$\circ$ & $32$ & $8$ & $8192$\\
total & $415$ & & $126464$\\
\hline
\end{tabular}
\end{table}

\section{Computer checks}\label{app:ver}

The ancillary scripts recheck printed data and run finite tests
(Table~\ref{tab:checks}). They use exact rational arithmetic and, for the
analytic bounds, Arb ball arithmetic \cite{Arb} through
\texttt{python-flint}~0.9.0 under Python~3.12.3.

\begin{center}\small\renewcommand{\arraystretch}{1.12}
\begin{minipage}{\textwidth}\centering
\captionof{table}{Computer checks.}\label{tab:checks}
\begin{tabular}{@{}>{\raggedright\arraybackslash}p{0.2\textwidth}>{\raggedright\arraybackslash}p{0.75\textwidth}@{}}
\hline
check & statements checked\\\hline
determinant&the exact seed determinant:
\eqref{eq:g4-polynomial} and Table~\ref{tab:g4-coefficients}. Finite tests:
$C_{ij}=3^i$ times the right side of \eqref{eq:laurent-moments} with
$d=j-i$, for $0\le i\le16$ and $|j-i|\le16$; the four
determinant evaluations in Lemma~\ref{lem:integral-identity} and their
product for $m\le3$ and three choices of $(s,k,h)$.\\
arithmetic&recomputes the arrangement of
\eqref{eq:subdivision-lines} for $R=\Rseam$ in exact rational arithmetic;
checks the line families of Table~\ref{tab:line-families}, every row of
Table~\ref{tab:slabs} (by evaluating $\mathfrak c$ from
\eqref{eq:residue-sets}--\eqref{eq:residue-rank} at two interior points),
\eqref{eq:inner-integral}, $A$ in \eqref{eq:tau}, $45/32$ and $-5/48$; rational
bounds for $\tau$ and $\alpha_0-\tau$; integrality of $K_mG_m$ for small $n$.\\
analytic&Lemmas~\ref{lem:reference-bound} and~\ref{lem:leaves}: the
chains, masses, conjugation symmetry and $|c|<r$ of
Tables~\ref{tab:reference-intervals} and~\ref{tab:reference-poisson};
Tables~\ref{tab:energy}, \ref{tab:caps} and~\ref{tab:depths},
\eqref{eq:Q-lower}, \eqref{eq:caps} and $\alpha_0$, with the bound of the
proof of Lemma~\ref{lem:reference-bound} on each of the $\nleaves$ leaves.\\
mod $17$&\eqref{eq:g4-mod-matrix}, \eqref{eq:g-mod-17}, and
\eqref{eq:block-congruence} for $t=1$ at $W=0,1,2$.\\\hline
\end{tabular}
\end{minipage}
\end{center}

\section{Proof of Fact~\ref{fact:energy}}\label{app:energy}

For $t>0$ let $g_t(y)=(\pi t)^{-1}e^{-|y|^2/t}$ on $\mathbb C$, a centred
Gaussian density. Then $g_s*g_t=g_{s+t}$ and $g_t(-y)=g_t(y)$. For $a,b>0$,
$\int_0^\infty(e^{-a/t}-e^{-b/t})\,dt/t=\log(b/a)$, and the integrand has
constant sign. With $a=|y|^2$ and $b=1$ this gives, for $y\ne0$,
\[
 -\log|y|=\frac\pi2\int_0^\infty\bigl(g_t(y)-g_t(1)\bigr)\,dt,\qquad
 \int_0^\infty|g_t(y)-g_t(1)|\,dt=\frac2\pi\bigl|\log|y|\bigr|.
\]
The hypothesis of Fact~\ref{fact:energy} makes the diagonal
$|\sigma|\otimes|\sigma|$-null and, by the second identity,
$g_t(u-w)-g_t(1)$ integrable for $|\sigma|\otimes|\sigma|\otimes dt$. By
Fubini's theorem and $\sigma(\mathbb C)=0$ (for each $t>0$, $g_t$ is
bounded, so $\iint g_t(1)\,d\sigma\,d\sigma=g_t(1)\sigma(\mathbb C)^2=0$),
\[
 -\iint\log|u-w|\,d\sigma(u)\,d\sigma(w)
 =\frac\pi2\int_0^\infty\iint g_t(u-w)\,d\sigma(u)\,d\sigma(w)\,dt.
\]
For fixed $t$, $g_t(u-w)=\int g_{t/2}(x-u)g_{t/2}(x-w)\,dx$ ($dx$ the
Lebesgue measure on $\mathbb C$), and the
corresponding triple integral against $|\sigma|\otimes|\sigma|$ is at most
$|\sigma|(\mathbb C)^2/(\pi t)$. So, by Fubini again,
\[
 \iint g_t(u-w)\,d\sigma(u)\,d\sigma(w)
 =\int_{\mathbb C}\Bigl(\int g_{t/2}(x-u)\,d\sigma(u)\Bigr)^2dx\ge0.
\]
Hence $\iint\log|u-w|\,d\sigma(u)\,d\sigma(w)\le0$. Compact support is not
used.

\section*{Use of AI tools}

The author devised the method of proof. Codex (OpenAI, GPT-6.1 Sol) computed the
reference measure $\nu$ and the other printed data, and wrote the Python
scripts in \texttt{anc/}. Claude Code (Anthropic,
Claude Opus 5.5) checked the scripts, the
data and the tables. The author checked every step
of the proof, every table and every reference, and takes full responsibility
for the content of the paper.

\section*{Statements and declarations}

\textbf{Competing interests.} The author has no competing interests.

\textbf{Funding.} No funding was received for this work.

{\rightskip=0pt plus 3em\relax
\textbf{Data availability.} All data used in the proof are printed in
Appendices~\ref{app:arithmetic}--\ref{app:reference-data}. The \texttt{anc/}
attachment holds the scripts of Appendix~\ref{app:ver} and a
\texttt{README.md}. The entry points are \texttt{check\_determinant.py},
\texttt{check\_arith.py}, \texttt{check\_analytic.py} and
\texttt{check\_mod17.py}. The helpers are \texttt{arith.py}, \texttt{caps.py},
\texttt{combine.py}, \texttt{energy.py}, \texttt{indep\_caps\_params.py},
\texttt{mod17\_replay.py}, \texttt{parse\_tex.py}, \texttt{refmeasure.py},
\texttt{tex2refjson.py} and \texttt{tex\_tables.py}. The scripts read the
tables and constants from the \TeX\ source of this manuscript. The analytic
replay writes to a temporary directory. The reports in
\texttt{anc/verification/} record the SHA-256 of the \TeX\ source they
checked. The file \texttt{anc/SCRIPTS\_SHA256SUMS} lists the SHA-256 of the
scripts and the \texttt{README.md}. Its own SHA-256 is\par}
\begin{center}\small
\texttt{890192149e125730cc0b9b0f85ddd7dab8e9ac391a6d14703849fb74d54f8cd7}.
\end{center}

\begin{thebibliography}{10}

\bibitem{Baker}
A.~Baker, \emph{Transcendental Number Theory}, 2nd ed., Cambridge Mathematical
Library, Cambridge University Press, Cambridge, 1990. MR1074572

\bibitem{CDT}
F.~Calegari, V.~Dimitrov, and Y.~Tang, \emph{The linear independence of $1$,
$\zeta(2)$, and $L(2,\chi_{-3})$}, preprint (2024), arXiv:2408.15403v2.

\bibitem{HK}
G.-N. Han and C.~Krattenthaler, \emph{Rectangular Scott-type permanents},
S\'eminaire Lotharingien de Combinatoire \textbf{43} (1999), Article B43g,
25 pp., arXiv:math/0003072. MR1741888

\bibitem{HW}
G.~H. Hardy and E.~M. Wright, \emph{An Introduction to the Theory of
Numbers}, 6th ed., revised by D.~R. Heath-Brown and J.~H. Silverman,
Oxford University Press, Oxford, 2008. MR2445243

\bibitem{Ha}
M.~Hata, \emph{Rational approximations to the dilogarithm}, Trans. Amer. Math.
Soc. \textbf{336} (1993), no.~1, 363--387,
\url{https://doi.org/10.1090/S0002-9947-1993-1147401-5}. MR1147401

\bibitem{Arb}
F.~Johansson, \emph{Arb: efficient arbitrary-precision midpoint-radius
interval arithmetic}, IEEE Trans. Comput. \textbf{66} (2017), no.~8,
1281--1292, \url{https://doi.org/10.1109/TC.2017.2690633}. MR3681746

\bibitem{Krat}
C.~Krattenthaler, \emph{Advanced determinant calculus: a complement},
Linear Algebra Appl. \textbf{411} (2005), 68--166,
\url{https://doi.org/10.1016/j.laa.2005.06.042}. MR2178686

\bibitem{DLMF}
\emph{NIST Digital Library of Mathematical Functions},
\url{https://dlmf.nist.gov/}, Release 1.2.8 of 2026-09-15, F.~W.~J. Olver
et al., eds.

\bibitem{Pr}
T.~Prellberg, \emph{Three irrationality results for the dilogarithm:
Rhin--Viola and Viola--Zudilin constructions at $-1/4$, $1/5$, and $-1/3$},
preprint (2026), arXiv:2609.06895v1.

\bibitem{RV05}
G.~Rhin and C.~Viola, \emph{The permutation group method for the dilogarithm},
Ann. Sc. Norm. Super. Pisa Cl. Sci. (5) \textbf{4} (2005), no.~3, 389--437,
\url{https://doi.org/10.2422/2036-2145.2005.3.02}. MR2185863

\bibitem{RV19}
G.~Rhin and C.~Viola, \emph{Linear independence of $1$, $\mathrm{Li}_1$ and
$\mathrm{Li}_2$}, Mosc. J. Comb. Number Theory \textbf{8} (2019), no.~1,
81--96,
\url{https://doi.org/10.2140/moscow.2019.8.81}. MR3864310

\bibitem{ST}
E.~B. Saff and V.~Totik, \emph{Logarithmic Potentials with External Fields},
Grundlehren der Mathematischen Wissenschaften, vol.~316, 1st ed., Springer, 1997,
\url{https://doi.org/10.1007/978-3-662-03329-6}. MR1485778

\bibitem{VZ}
C.~Viola and W.~Zudilin, \emph{Linear independence of dilogarithmic values},
J. Reine Angew. Math. \textbf{736} (2018), 193--223,
\url{https://doi.org/10.1515/crelle-2015-0030}; preprint version:
Max-Planck-Institut f\"ur Mathematik Preprint Series 2014 (26),
\url{https://archive.mpim-bonn.mpg.de/id/eprint/1914/}. MR3769989

\end{thebibliography}

\end{document}
